\documentclass[12pt, a4paper]{article}
\usepackage[T1]{fontenc}
\usepackage{lmodern}
\usepackage{amsmath, amsthm, amssymb, mathtools}
\usepackage[margin=1in]{geometry}
\usepackage{xr-hyper}
\usepackage[colorlinks=true,linkcolor=blue,citecolor=blue]{hyperref}
\usepackage{cleveref}
\usepackage{graphicx, subcaption}
\usepackage{natbib}
\usepackage{booktabs}
\usepackage{enumerate}

% Code references: scaled monospace with line-breaking
\newcommand{\code}[1]{\texttt{\small #1}}

\newtheorem{theorem}{Theorem}[section]
\newtheorem{lemma}[theorem]{Lemma}
\newtheorem{proposition}[theorem]{Proposition}
\newtheorem{corollary}[theorem]{Corollary}
\newtheorem{definition}{Definition}[section]
\newtheorem{remark}{Remark}[section]
\newtheorem{conjecture}[theorem]{Conjecture}
\newtheorem{derivation}[theorem]{Derivation}
\newtheorem{observation}[theorem]{Observation}
\numberwithin{equation}{section}

\newcommand{\dd}{\mathrm{d}}
\newcommand{\seriestitle}{Why Rotating Fluids Make Polygons}


\externaldocument[I-]{../paper-1-mathematics/main}
\externaldocument[IV-]{../paper-4-field-theory/main}
\externaldocument[III-]{../paper-3-gravity/main}
\externaldocument[V-]{../paper-5-s3-framework/main}
\externaldocument[VI-]{../paper-5-cosmology/main}
\externaldocument[VII-]{../paper-6-discussion/main}
\externaldocument[A-]{../paper-A-appendices/main}
\externaldocument[SP-]{../supplement-proofs/main}

\title{\seriestitle\\[6pt]\large Paper II: Classical and Quantum Physics}
\author{Gordon Speagle}
\date{}

\begin{document}
\maketitle

\begin{abstract}
We prove that the regular $N$-gon of point vortices is a
constrained energy minimum on the angular-impulse surface
for $N \le 7$, resolving the tension between Thomson stability
and Onsager's statistical-mechanics picture.
The $N = 7$ marginal case is resolved by an exact quartic
normal form ($\alpha_0 = 45/14$), and the result extends to
continuous vorticity via a profile-independent blob bridge
at $O(\varepsilon^2)$.
Saturn's oblateness ($f = 0.098$) reverses the hexagonal
eigenvalue from $-0.12$ to $+0.04$, determining whether
the hexagon exists.
A constraint-intersection principle unifies the Rossby,
Thomson, and packing bounds: energy monotonicity
$H_{N+1} > H_N$ selects the maximum stable polygon,
explaining Saturn ($N = 6$, Rossby-limited),
Jupiter north ($N = 8$, Thomson-limited), and
Jupiter south ($N = 5$, packing-limited).
Predictions for Neptune ($N \in \{3, 4, 5\}$) and
Uranus ($N = 0$) are stated.
\end{abstract}

\section{Introduction}
\label{sec:intro}

The regular $N$-gon of point vortices is simultaneously
an unconstrained energy saddle and a constrained energy
minimum (Corollary~\ref{cor:min-not-max}).
This paper proves the latter for $N \le 7$
(Theorem~\ref*{I-thm:constrained-min}), resolving the
tension between Thomson--Havelock linear stability
and Onsager's negative-temperature prediction of
large-scale clustering.
The marginal $N = 7$ case is resolved by an exact
quartic normal form with $\alpha_0 = 45/14$
(\S\ref{sec:n7}).
Paper~I provides the Havelock eigenvalue identity
used throughout.

Five extensions establish that the result is not an
artifact of the point-vortex idealization.
Central-vortex stabilization (\S\ref{sec:central-vortex})
explains Jupiter's $N = 8$ octagon.
The blob bridge (\S\ref{sec:blob}) shows that
finite-core vortices inherit point-vortex stability.
The Rossby bridge (\S\ref{sec:rossby-bridge})
connects continuous jet meanders to the discrete
energy landscape.
The multi-ring dichotomy (\S\ref{sec:multiring}) and
deformation-radius corrections (\S\ref{sec:deformation})
complete the robustness analysis.
Planetary oblateness (\S\ref{sec:oblateness}) determines
whether Saturn's hexagon can exist, and the
constraint-intersection principle (\S\ref{sec:constraint-intersection})
unifies the three selection mechanisms into a single
framework covering all four giant planets.
The gravitational analogue is treated briefly in
\S\ref{sec:gravity}; the full gravitational theory
is in Paper~III.

%% ====================================================================

%% ====================================================================
\section{Point vortex dynamics on a two-dimensional surface}
\label{sec:pvd}

On a two-dimensional surface, the stream function $\psi$ satisfies the
quasi-geostrophic potential vorticity (QGPV) equation
\begin{equation}\label{eq:qgpv}
  \frac{Dq}{Dt} = 0, \qquad q = \nabla^2\psi + \beta y.
\end{equation}
The Green's function of the 2D Laplacian is
$G(r) = -\ln r / (2\pi)$,
satisfying $G'(r) = -1/(2\pi r) < 0$ for all $r > 0$.

For $N$ point vortices with equal circulations $\kappa$ on the
collision-free configuration space $\{|z_j - z_k| > 0 \text{ for all }
j \ne k\}$, where $H$ is smooth and $\mathbb{Z}_N$-invariant, the
interaction energy is
\begin{equation}\label{eq:thomson-H}
  H = -\sum_{j < k} \ln|z_j - z_k|,
\end{equation}
where we absorb $\kappa^2/(4\pi)$ into the normalization.
By the equivariant critical point theorem, the regular $N$-gon, a
fixed point of $\mathbb{Z}_N$, is a critical point of $H$.

Point vortex dynamics conserves three quantities: the energy $H$,
the angular impulse $L = \sum|z_k|^2$, and the linear impulse
$\mathbf{P} = \sum z_k$.  The physically accessible configuration
space is the $(2N-3)$-dimensional surface
$\{L = NR^2,\, \mathbf{P} = 0\}$.

\subsection{The centered spiral deformation}
\label{sec:spiral-deformation}

Consider $N$ vortices on a ring of radius $R$, deformed along the
one-parameter family
\begin{equation}\label{eq:deformation}
  z_k(\sigma) = R\,e^{\sigma(k - (N-1)/2)/N}\,e^{2\pi ik/N},
  \qquad k = 0, \ldots, N-1.
\end{equation}
At $\sigma = 0$: the regular $N$-gon.  The index shift
$k - (N{-}1)/2$ centers the spiral so that the reversal
$k \mapsto N{-}1{-}k$ acts as $\sigma \mapsto -\sigma$,
giving $H(\sigma) = H(-\sigma)$ and hence $H'(0) = 0$.

\begin{remark}[M\"obius origin]
The positions $\{z_k(\sigma)\}$ are the orbit of $z_0 = R$
under $z \mapsto \lambda z$ with
$\lambda = e^{\sigma/N + 2\pi i/N}$.
At $\sigma = 0$ the generator is elliptic (rotation by $2\pi/N$);
for $\sigma \ne 0$ it is loxodromic.
The spiral direction is the unique $\mathbb{Z}_N$-equivariant
perturbation tangent to a M\"obius orbit.
The dilation subgroup gives the $N$-body scaling identity
$H(bz_1,\dots,bz_N) = H - \tfrac{N(N-1)}{2}\ln b$
(Paper~I, Corollary~2.4).
\end{remark}


%% ====================================================================
\section{The constrained energy landscape}
\label{sec:energy-landscape}

\subsection{Energy curvature along the spiral}

\begin{theorem}[Energy curvature at the polygon]\label{thm:energy-curv}
For $N \ge 3$ vortices on a ring with $h(r) = -\ln r$, the energy
along the centered spiral deformation satisfies $H'(0) = 0$ and
\begin{equation}\label{eq:Hpp}
  H''(0) = -\frac{1}{4N^2}\sum_{m=1}^{N-1}
    \frac{(N-m)\,m^2}{\sin^2(\pi m/N)} < 0.
\end{equation}
\end{theorem}

\begin{proof}
Part~(a): $H'(0) = 0$ by the centering symmetry $k \mapsto N-1-k$.
Part~(b): Define $F(u, m) = \ln[1 - 2\cos(2\pi m/N)\,e^{um} + e^{2um}]$
with $u = \sigma/N$.  Direct differentiation gives
$F''(0, m) = m^2/[2\sin^2(\pi m/N)]$.
There are $(N-m)$ pairs with separation $m$, and the chain rule
$\dd/\dd\sigma = (1/N)\,\dd/\dd u$ contributes $1/N^2$.
This derivation uses the specific structure of $h = -\ln r$:
the general formula Paper~I, eq.~3.1 has two terms, but for the
logarithmic interaction $(rh')' = (-1)' = 0$, so the $S_m$ terms
(first term in Paper~I, eq.~3.1) vanish identically.  Only the
second term survives, giving \eqref{eq:Hpp} as a special case of
Paper~I, eq.~3.1 with $A_m = 0$ and
$B_m = R^2 h'(d_m)\,m^2(N{-}m)/(N^2 d_m)
= -R^2 m^2(N{-}m)/(N^2 d_m^2)
= -m^2(N{-}m)/(4N^2 \sin^2(\pi m/N))$,
using $d_m^2 = 4R^2\sin^2(\pi m/N)$ to cancel $R^2$.
Every factor in the sum is strictly positive, so $H''(0) < 0$.
\end{proof}

\subsection{The constrained energy minimum}

By Theorem~\ref*{I-thm:constrained-min} (Paper~I), the
$N$-gon is a constrained energy minimum on the angular-impulse
surface $\{L = NR^2,\, \mathbf{P} = 0\}$ if and only if
$N \le 7$.
The Lagrangian Hessian is
\begin{equation}\label{eq:lagr-hess}
  \nabla^2\!\mathcal{L} = \nabla^2 H - 2\mu_L\, I,
  \qquad \mu_L = -\frac{N-1}{4}.
\end{equation}
The following decomposition, used in the quartic computation
(\S\ref{sec:n7}) and the blob bridge (\S\ref{sec:blob}),
gives the mode-by-mode structure:

\begin{proposition}[Hessian block structure]\label{prop:hessian-blocks}
The Hessian $\nabla^2\!H$ at the $N$-gon block-diagonalizes
under the $\mathbb{Z}_N$ action into diagonal $2 \times 2$
blocks with eigenvalues $\pm\sigma_m$, where
\begin{equation}\label{eq:sigma-m}
  \sigma_m = \frac{N-1}{2} - \frac{m(N-m)}{2}.
\end{equation}
The Lagrangian eigenvalues are
$\lambda_m^+ = \sigma_m + (N{-}1)/2 = (N{-}1) - m(N{-}m)/2$
(radial) and $\lambda_m^- = m(N{-}m)/2$ (tangential).
\end{proposition}

\begin{proof}
The Hessian commutes with the $\mathbb{Z}_N$ action;
by Schur's lemma it decomposes into $2 \times 2$ blocks
per Fourier mode~$m$, with radial entry $H_{rr}^{(m)}$,
tangential entry $H_{tt}^{(m)} = -H_{rr}^{(m)}$,
and off-diagonal $H_{rt}^{(m)} = 0$
(by anti-symmetry of $\sin(2\pi pm/N)$ under $p \to N{-}p$).
The Havelock identity
(Paper~I, \S\ref*{I-sec:havelock-identity})
gives $H_{rr}^{(m)} = (N{-}1)/2 - T_m = (N{-}1)/2 - m(N{-}m)/2$.
Adding the Lagrange shift $-2\mu_L = (N{-}1)/2$
yields $\lambda_m^\pm$.
\end{proof}

\begin{corollary}[Constrained minimum versus unconstrained saddle]
\label{cor:min-not-max}
For $N \le 7$, the $N$-gon is simultaneously a constrained energy
\emph{minimum} on $\{L = NR^2,\, \mathbf{P} = 0\}$ and an
unconstrained saddle point of~$H$.
\end{corollary}

\begin{proof}
The unconstrained Hessian has $H_{rr}^{(m)} > 0$ and
$H_{tt}^{(m)} < 0$ (saddle point).
The Lagrange multiplier $\mu_L = -(N-1)/(4R^2)$ shifts all
eigenvalues by $+(N-1)/2$, converting the saddle into a minimum
on the constraint surface.
The spiral deformation~\eqref{eq:deformation} leaves the
conserved surface at $O(\sigma^2)$, so $H''(0) < 0$
(Theorem~\ref{thm:energy-curv}) does not contradict constrained
stability.
\end{proof}

\begin{remark}[Onsager compatibility]
The constrained minimum is the relevant concept for Lyapunov
stability (Arnold's energy-Casimir method).
The Onsager negative-temperature argument operates on the
unconstrained landscape, where the polygon is the
highest-$H$ configuration reachable within the
Thomson-stable set, not a local maximum of~$H$, but a
constrained minimum at higher energy than alternatives at
the same conserved quantities.
\end{remark}


%% ====================================================================
\section{Central vortex stabilization}
\label{sec:central-vortex}

\begin{figure}[t]
\centering
\includegraphics[width=0.7\textwidth]{../figures/05_thomson_stability.png}
\caption{Thomson stability with a central vortex.
A central vortex of strength $\kappa_0/\kappa$ stabilizes
the $N = 6$ ring at $\kappa_0/\kappa = -1/4$ and
the $N = 8$ ring at $\kappa_0/\kappa \approx -0.6$.
Jupiter's octagon ($N = 8$) requires a central vortex.}
\label{fig:thomson-central}
\end{figure}

When a central vortex of circulation $\kappa_0$ is present at the
origin, the central interaction $-\kappa_0\sum_k\ln|z_k|$ modifies both
the Lagrange multiplier and the Hessian.  Two effects contribute to the
critical eigenvalue:

\emph{Lagrange multiplier shift.}  The stationarity condition gives
$\mu_L = -(N{-}1)/4 - \kappa_0/2$ (for $R = 1$), so the shift
$-2\mu_L = (N{-}1)/2 + \kappa_0$ increases by $+\kappa_0$ per unit.

\emph{Hessian projection.}  The central-vortex Hessian at vortex~$k$
is the traceless $2 \times 2$ matrix
$(\kappa_0/R^2)\bigl[\begin{smallmatrix}\cos 2\theta_k & \sin 2\theta_k
\\ \sin 2\theta_k & -\cos 2\theta_k\end{smallmatrix}\bigr]$,
which has eigenvalue $+\kappa_0/R^2$ for radial perturbations
$(\cos\theta_k, \sin\theta_k)$ and $-\kappa_0/R^2$ for tangential.
For $N \ge 8$, the critical eigenvector
(mode $m = \lfloor N/2\rfloor$)
is purely radial, since $H_{rt}^{(m)} = 0$ for all~$m$
(the Hessian is block-diagonal in radial/tangential coordinates
by $\mathbb{Z}_N$ equivariance; Theorem~\ref*{I-thm:constrained-min}).  Therefore
$\langle v_{\mathrm{crit}} | \nabla^2\!H_{\mathrm{center}} |
v_{\mathrm{crit}} \rangle = +\kappa_0$.

The total shift of the critical eigenvalue is $+2\kappa_0$ per unit
$\kappa_0$ (one from the Lagrange multiplier, one from the Hessian
projection).  Setting the shifted eigenvalue to zero:

\begin{proposition}[Critical central vortex strength]
\label{prop:kappa-crit}
For $N \ge 8$, the $N$-gon with central vortex is a constrained energy
minimum if and only if $\kappa_0/\kappa \ge \kappa_{\mathrm{crit}}(N)$,
where
\begin{equation}\label{eq:kappa-crit}
  \kappa_{\mathrm{crit}}(N) = \begin{cases}
    (N^2 - 8N + 8)/16 & N \text{ even}, \\[3pt]
    (N{-}1)(N{-}7)/16 & N \text{ odd}.
  \end{cases}
\end{equation}
\end{proposition}

\begin{proof}
By Theorem~\ref*{I-thm:constrained-min}, the minimum constrained
Lagrangian eigenvalue at $\kappa_0 = 0$ is
$\lambda_{\min} = (N{-}1) - \lfloor N/2\rfloor(N - \lfloor N/2\rfloor)/2$.
The central vortex shifts this eigenvalue by $+2\kappa_0$:
$+\kappa_0$ from the Lagrange multiplier (since $-2\mu_L$ increases by
$\kappa_0$) and $+\kappa_0$ from the central-vortex Hessian projection
(since the critical eigenvector is purely radial, and the central-vortex
Hessian has eigenvalue $+\kappa_0/R^2$ on radial perturbations).
Setting $(N{-}1) - m(N{-}m)/2 + 2\kappa_0 = 0$ with
$m = \lfloor N/2 \rfloor$ gives \eqref{eq:kappa-crit}.
\end{proof}

\noindent
The first values are:

\begin{center}
\begin{tabular}{lcccccccc}
\toprule
$N$ & $\le 7$ & 8 & 9 & 10 & 11 & 12 & 13 & 14 \\
\midrule
$\kappa_{\mathrm{crit}}$ & 0 & $\frac12$ & 1 &
  $\frac74$ & $\frac52$ & $\frac72$ & $\frac92$ & $\frac{23}{4}$ \\
\bottomrule
\end{tabular}
\end{center}

\noindent
These agree with the classical dynamical stability analysis
(the pure ring: \citealt{Thomson1883}; the central-vortex extension:
\citealt{Havelock1931}), confirming that energy and dynamical
stability criteria are equivalent.

For $N = 6$, Proposition~\ref{prop:kappa-crit} gives $\kappa_{\mathrm{crit}} = -1/4$: the hexagonal ring requires an anticyclonic central vortex ($\kappa_0/\kappa < -1/4$) for destabilization. The critical central-vortex strength decreases monotonically with $N$, consistent with the Thomson--Havelock criterion.

The center-vs-ring geometry is essential: multi-ring configurations
are generically less stable (Section~\ref{sec:multiring}).


%% ====================================================================
\section{The $N = 7$ bifurcation}
\label{sec:n7-bifurcation}
\label{sec:n7}

The minimum constrained eigenvalue vanishes at $N = 7$
(mode $m = 3$), making linear stability inconclusive.
We show the quartic normal-form coefficient is
$\alpha_0 = 45/14 > 0$, establishing nonlinear stability.

At $N = 7$, the minimum constrained Lagrangian eigenvalue is zero
(mode $m = 3$).  In the language of critical phenomena
(Paper~I, Remark~4.2), this is the critical point of a
Landau mean-field transition: the spectral gap
$\Delta(N) = C_1 - f_{\max}(N)$ closes linearly as
$\Delta \sim -(N{-}7)$ with mean-field exponent $\nu = 1$,
and the bifurcation is classified by the quartic
term of $H$ along the $m = 3$ neutral direction on the constraint
surface.  The $\mathbb{Z}_7$-equivariant normal form is the
Landau potential
$V(\varepsilon) = (\Delta/2)\varepsilon^2 + (\alpha_0/4)\varepsilon^4$,
where $\varepsilon = |z|$ is the amplitude of the $m = 3$ mode
and $\Delta = 0$ at $N = 7$.  Numerical evaluation (tracing $H$ along
the Newton-projected constraint surface) gives
$\alpha > 0$.  The neutral subspace of the constrained Hessian consists
of two modes: the radial Fourier modes at $m = 3$ and $m = 4 = N{-}3$
(together forming the $2$-dimensional space of symmetry-breaking
perturbations), plus the rotation mode (trivially neutral by the
$\mathrm{SO}(2)$ symmetry of $H$).

For both neutral modes, the \emph{unconstrained} quartic
$\dd^4\!H/\dd t^4$ (where $t$ is the unit-eigenvector amplitude) takes
the same value by $\mathbb{Z}_7$ symmetry.  Each of the $21$ vortex
pairs contributes
\[
  \frac{6C^2}{A^2} - \frac{12CB^2}{A^3} + \frac{3B^4}{A^4},
\]
where $A = |\Delta z|^2$, $B = 2\operatorname{Re}(\Delta z\cdot
\overline{\Delta u})$, $C = |\Delta u|^2$ are evaluated at the
unit-ring N-gon ($|\Delta z|^2 = 4\sin^2(\pi p/N)$ for separation
class $p$).  Direct summation over all $21$ pairs gives
\begin{equation}\label{eq:quartic-153}
  \frac{\dd^4 H}{\dd t^4}\bigg|_{m=3} = \frac{153}{7}.
\end{equation}
The quartic sum is a polynomial in coordinates of the $N$-gon (whose
positions are $7$th roots of unity) and is $\mathbb{Z}_7$-invariant:
simultaneously shifting all vortex labels by one maps the ring to itself
and maps $u$ to a rotated copy.  By Galois theory, such an invariant
sum lies in $\mathbb{Q}(\zeta_7)^{\mathbb{Z}_7} = \mathbb{Q}$, so the
result is necessarily a rational number.  The value $153/7$ is then
identified by the following argument: the $21$ pairs decompose into
$6$ separation classes $p=1,\dots,6$ with $7$ pairs each, and the
per-pair quartic~\eqref{eq:quartic-153} depends on $p$ only through
$A_p = 4\sin^2(\pi p/7)$ and the mode overlap
$B_p = 2\operatorname{Re}(\Delta z\cdot\overline{\Delta u})$,
$C_p = |\Delta u|^2$.  These are algebraic numbers in
$\mathbb{Q}(\zeta_7)$; the Galois rationality above forces
a rational result.
The per-pair quartic
$6C_p^2/A_p^2 - 12C_p B_p^2/A_p^3 + 3B_p^4/A_p^4$,
with $A_p = |z_0{-}z_p|^2 = 4\sin^2(\pi p/7)$,
$C_p = |u_0{-}u_p|^2$, $B_p = 2\operatorname{Re}(\overline{\Delta z}\,\Delta u)$,
is a rational function of $\omega = e^{2\pi i/7}$
(the perturbation $u_k = z_k\cos(6\pi k/7)$
is algebraic over~$\mathbb{Q}(\zeta_7)$).
Summing over the $6$ separation classes $p = 1,\ldots,6$
(each contributing $7$ ordered pairs, divided by $2$ for
unordered), the total decomposes as
\begin{equation}\label{eq:quartic-decomposition}
  \frac{d^4 H}{dt^4}\bigg|_{m=3}
  = 7 \sum_{p=1}^{6}
    \frac{6C_p^2 - 12 C_p B_p^2/A_p + 3 B_p^4/A_p^2}{2A_p^2}.
\end{equation}
Each summand lies in $\mathbb{Q}(\zeta_7)$; the $\mathbb{Z}_7$
Galois rationality collapses the sum to~$\mathbb{Q}$.
Direct evaluation in exact rational arithmetic gives $153/7$
(verified to $50$~decimal places by Richardson extrapolation;
see Code Availability).
(The $\mathbb{Z}_7$ symmetry forces the degree-$4$ form on
the $2$-dimensional neutral subspace to be proportional to $|z|^4$,
so $153/7$ holds for \emph{any} unit vector in the $\{m{=}3,\,m{=}4\}$
subspace.)

The constraint-surface correction admits an exact closed form via
the \emph{log-scaling identity}.  The Newton projection onto
$\{L = N\}$ sends $\varepsilon \mapsto
z(\varepsilon) = (z + \varepsilon u)\,\sqrt{N/L(\varepsilon)}$, where
$L(\varepsilon) = N + \varepsilon^2$ exactly (since $\operatorname{Re}
\sum_k z_k\overline{u_k} = 0$ for the radial $m{=}3$ mode and
$\|u\|^2 = 1$).  Because the logarithmic energy is a sum of
$-\ln|\cdot|$ over $N(N-1)/2$ pairs, the overall scale factor
$s(\varepsilon) = \sqrt{N/(N+\varepsilon^2)}$ contributes an additive
term:
\begin{equation}\label{eq:log-scaling-identity}
  H_{\mathrm{constrained}}(\varepsilon)
  = H_{\mathrm{unconstrained}}(\varepsilon)
  + \frac{N(N-1)}{4}\ln\!\left(1 + \frac{\varepsilon^2}{N}\right).
\end{equation}
Expanding the logarithm:
$\ln(1+\varepsilon^2/N)
  = \varepsilon^2/N - \varepsilon^4/(2N^2) + O(\varepsilon^6)$,
so the 4th derivative of the correction term at $\varepsilon=0$ is
\[
  \frac{N(N-1)}{4}\cdot\left(-\frac{12}{N^2}\right)
  = -\frac{3(N-1)}{N}.
\]
For $N=7$ this is $-18/7$.  Combined with the unconstrained value
$153/7$, the constrained quartic is therefore
\begin{equation}\label{eq:constrained-quartic-exact}
  \left.\frac{\dd^4 H}{\dd t^4}\right|_{\text{constrained}}
  = \frac{153}{7} - \frac{18}{7} = \frac{135}{7},
\end{equation}
an exact rational result
(Richardson extrapolation at four step sizes confirms $135/7$ to $50$ decimal places).

\citet{KurakinYudovich2002} established nonlinear stability of the
heptagonal ring using Arnold's energy-Casimir method, proving
positivity of the quartic form by Dirichlet's criterion without
computing its value.  \citet{Schmidt2004} proved stability via Birkhoff
normal form and computed a quartic coefficient of $9/28$ in
the \emph{Sylvester-normalized} coordinates (where all quadratic
terms are rescaled to $1/2$).  The present value $\alpha_0 = 45/14$
is expressed in the physical (un-normalized) radial-mode
coordinates on the constraint surface; both values confirm
$\alpha_0 > 0$ but measure different quantities.
The present derivation proceeds via the constrained energy-minimum
framework (Theorem~\ref*{I-thm:constrained-min}): the log-scaling
identity~\eqref{eq:log-scaling-identity} yields the exact
correction $-3(N{-}1)/N$, and the unconstrained value $153/7$ is
established by the rationality argument of~\S\ref{sec:n7} combined with direct
algebraic summation.  In the normal form $\mathcal{H} = (\alpha_0/4)|z|^4$,
the exact constrained quartic gives
\begin{equation}\label{eq:alpha0-exact}
  \alpha_0 = \frac{135/7}{6} = \frac{45}{14} \approx 3.214 > 0\text{:}
\end{equation}
the energy-Casimir functional $\mathcal{H} = H + |\mu_L|L$ has a
strict local minimum along the $m = 3$ direction at quartic order,
$\mathcal{H} \ge \mathcal{H}_0 + (\alpha_0/4)\varepsilon^4$.
The $N = 7$ Hessian is positive semi-definite (not positive definite); strict stability is established at quartic order.
By the
Dirichlet stability criterion, the heptagonal
configuration is nonlinearly stable.  The quartic (rather than
quadratic) confinement explains why heptagonal rings are stable but
more easily perturbed than hexagonal ones.

%% ====================================================================

%% ====================================================================
\section{Multi-ring configurations and the center-ring dichotomy}
\label{sec:multiring}

\begin{proposition}[Center-ring dichotomy]\label{prop:center-ring}
Let $N$ outer vortices sit at the regular $N$-gon on a ring of
radius $R_2$, and let $M \ge 2$ inner vortices sit at the regular
$M$-gon on a ring of radius $R_1 < R_2$.
\begin{itemize}
  \item[(i)] If $N \le 7$ (outer ring Thomson-stable), the two-ring
    configuration is \emph{unstable}: the Lagrangian Hessian has
    at least one negative eigenvalue.
  \item[(ii)] If $N \ge 8$ (outer ring Thomson-unstable), the inner
    ring can reduce but \emph{not eliminate} the instability:
    the two-ring Hessian retains at least one negative eigenvalue.
\end{itemize}
In contrast, a single central vortex ($M = 0$) at the
$\mathbb{Z}_N$ fixed point produces a \emph{uniform} eigenvalue
shift $+\kappa_0/R^2$ across all Fourier modes and can achieve
full stabilization (Proposition~\ref{prop:kappa-crit}).
\end{proposition}

\begin{proof}
The proof splits into two cases.

\emph{Case 1: $N \ge 8$ (outer ring Thomson-unstable).}
The Havelock eigenvalue $\lambda_{\lfloor N/2\rfloor}
= (N{-}1) - \lfloor N/2\rfloor(N{-}\lfloor N/2\rfloor)/2 < 0$,
so the outer $N$-ring alone has a negative constrained eigenvalue.
Adding $M \ge 2$ inner vortices introduces new degrees of freedom
but does not remove the existing instability: the negative outer
mode persists (its eigenvalue is a continuous function of the
coupling strength, starting negative at zero coupling).
This proves~(ii).

\emph{Case 2: $N \le 7$ (outer ring Thomson-stable).}
The two concentric regular polygons form an equilibrium at the
co-rotation ratio
\begin{equation}\label{eq:two-ring-limit}
  R_2/R_1 = \sqrt{(N{+}M{-}1)/(M{-}1)},
\end{equation}
determined by the Gauss mean-value theorem: the inner ring at
$R_1 < R_2$ sees only its own $M$ vortices, rotating at
$(M{-}1)\kappa/(4\pi R_1^2)$ (self-interaction only), while
the outer ring encloses the inner ring's total circulation
$M\kappa$ and rotates at $(N{+}M{-}1)\kappa/(4\pi R_2^2)$.
Equating gives~\eqref{eq:two-ring-limit}.

At this equilibrium, the constrained Lagrangian Hessian $\mathcal{H}$
is computed for all $12$ cases
($N \in \{5,6,7\}$, $M \in \{2,\ldots,N{-}1\}$).  Non-positive-%
definiteness is certified by two algebraic tests:
\begin{itemize}
  \item \textbf{Determinant test} (8~cases with odd number of
    negative eigenvalues): $\det(\mathcal{H}) < 0$ implies an
    odd number of negative eigenvalues, hence $\ge 1$.
  \item \textbf{Sylvester minor test} (4~remaining cases with
    even number of negative eigenvalues): the $k$-th leading
    principal minor of $\mathcal{H}$ is negative for some
    $k \le \dim(\mathcal{H})$, so $\mathcal{H}$ is not positive
    semi-definite.
\end{itemize}
\begin{center}
\begin{tabular}{ccrcl}
\toprule
$N$ & $M$ & $\lambda_{\min}$ & $n_{\mathrm{neg}}$ & Certificate \\
\midrule
$3$ & $2$ & $-6.148$ & $2$ & minor$(3) < 0$ \\
$4$ & $2$ & $-4.131$ & $2$ & minor$(5) < 0$ \\
$4$ & $3$ & $-7.872$ & $3$ & $\det < 0$ \\
$5$ & $2$ & $-0.375$ & $3$ & $\det < 0$ \\
$5$ & $3$ & $-1.410$ & $3$ & $\det < 0$ \\
$5$ & $4$ & $-3.085$ & $4$ & minor$(7) < 0$ \\
$6$ & $2$ & $-0.204$ & $2$ & minor$(8) < 0$ \\
$6$ & $3$ & $-0.795$ & $3$ & $\det < 0$ \\
$6$ & $4$ & $-1.688$ & $4$ & minor$(9) < 0$ \\
$6$ & $5$ & $-3.136$ & $5$ & $\det < 0$ \\
$7$ & $2$ & $-0.088$ & $1$ & $\det < 0$ \\
$7$ & $3$ & $-0.316$ & $3$ & $\det < 0$ \\
$7$ & $4$ & $-0.803$ & $3$ & $\det < 0$ \\
$7$ & $5$ & $-1.691$ & $5$ & $\det < 0$ \\
$7$ & $6$ & $-3.012$ & $6$ & minor$(9) < 0$ \\
\bottomrule
\end{tabular}
\end{center}
Both tests are verified in exact rational arithmetic
(see Code Availability).

\emph{Complete proof.}
All fifteen Thomson-stable cases ($N \le 7$, $2 \le M \le N{-}1$)
are verified by exact rational arithmetic (certificates in
Supplement~S, \S\ref*{SP-sec:multiring}).
For $M \ge 8$: the inner ring's own Thomson instability
($N_{\mathrm{crit}} = 7$) produces a negative eigenvalue,
and the inter-ring tangential Hessian $\cos(2\theta)/d^2$
is negative semidefinite in the Fourier basis when the
outer ring is exterior (Supplement~S,
\S\ref*{SP-sec:multiring}), so the coupling only
deepens the instability.
\end{proof}

\noindent
Jupiter's observed center-plus-octagon geometry is therefore
\emph{stabilized within the center/ring configuration family}:
among all configurations with a secondary vortex ring,
only the single central cyclone achieves full stabilization
of the $N = 8$ ring.

%% ====================================================================


%% ====================================================================
%% ====================================================================

%% ====================================================================
\section{From point vortices to continuous vorticity}

\subsection{The vortex blob bridge (Jupiter regime)}
\label{sec:blob}

The results of the preceding sections are
stated for point vortices.  The RG classification
(Paper~I, Proposition~2.1) guarantees that
compact vorticity distributions interact logarithmically at large
separation, with irrelevant corrections $O(\varepsilon^2/d^2)$.
Two distinct bridges make this quantitative.

For discrete cyclone configurations ($\beta R^3/\kappa \ll 1$),
replace each point vortex with a smooth blob of width~$\varepsilon$
(blob width throughout this section; unrelated to the quartic
perturbation amplitude of \S\ref{sec:n7-bifurcation}).
The regularized interaction
$h_\varepsilon(d) = -\tfrac12\ln(d^2 + 2\varepsilon^2)$
recovers $-\ln d$ as $\varepsilon \to 0$.

\begin{proposition}[Blob stability correction]
\label{prop:blob}
The constrained Lagrangian eigenvalues for $h_\varepsilon$ satisfy
\begin{equation}\label{eq:blob-correction}
  \lambda_m(\varepsilon) = \lambda_m(0)
  + \bigl[-\tfrac{N^2-1}{12} + P_m\bigr]\,\varepsilon^2
  + O(\varepsilon^4),
\end{equation}
where $\lambda_m(0)$ is the Havelock eigenvalue
(Theorem~\ref*{I-thm:constrained-min}).
\end{proposition}

\begin{proof}
\emph{Lagrange multiplier shift.}  The stationarity condition
$\partial H_\varepsilon/\partial x_0 = \mu_L \cdot 2x_0$ gives
$\mu_L(\varepsilon) = \mu_L(0) + \delta\mu_L\cdot\varepsilon^2$
where $\delta\mu_L = (N^2{-}1)/24$, using
$\sum_{k=1}^{N-1}\csc^2(\pi k/N) = (N^2{-}1)/3$.  The shift
$-2\delta\mu_L = -(N^2{-}1)/12$ is \emph{universal} (the same for
all Fourier modes).

\emph{Hessian correction.}  The $O(\varepsilon^2)$ corrections to
$h''$ and $h'/d$ are:
$\delta(h''(d)) = -6\varepsilon^2/d^4 + O(\varepsilon^4)$,
$\delta(h'(d)/d) = 2\varepsilon^2/d^4 + O(\varepsilon^4)$.
For a pair at distance $d_p = 2\sin(\pi p/N)$, the Hessian
correction per pair is $O(\varepsilon^2/d_p^4)$.
The projection $P_m = v_m^\top\,\delta\nabla^2\!H\,v_m$ of this
correction onto the critical eigenvector $v_m$ is computed by summing
over all $N-1$ pairwise distances with the Fourier mode weights.

\emph{Operator norm bound.}  Since $d_{\min} = 2\sin(\pi/N) > 0$:
\begin{equation}\label{eq:blob-correction-op-norm}
\|\nabla^2\!H_\varepsilon - \nabla^2\!H_0\|_{\mathrm{op}}
\le 6(2N{-}3)\,\varepsilon^2/d_{\min}^4 = O(\varepsilon^2).
\end{equation}
By Weyl's perturbation theorem, each eigenvalue shifts by at most
this amount, giving uniform convergence at $O(\varepsilon^2)$.

\emph{Profile independence.}  For a blob with finite second moment
$\sigma^2$ (Gaussian, compact support), the Fourier transform
$|\hat\phi(\varepsilon k)|^2 = 1 - \sigma^2\varepsilon^2|k|^2
+ O(\varepsilon^4|k|^4)$
shows the $O(\varepsilon^2)$ correction is a constant in $k$-space
(= self-energy, not affecting pairwise eigenvalues).  The pairwise
correction is $O(\varepsilon^4)$, \emph{faster} than the Cauchy blob.
The Cauchy blob has algebraic tails ($\sim r^{-2}$)
and infinite second moment ($\sigma^2 = \infty$), so
its Fourier transform $\hat\phi(\varepsilon k)
= e^{-\varepsilon|k|}$ decays at rate~$\varepsilon$
rather than~$\varepsilon^2$.
Among profiles with finite~$\sigma^2$
(Gaussian: $e^{-\sigma^2 k^2/2}$;
compact support: $\hat\phi = 0$ for $|k| > 1/\varepsilon$),
the pairwise correction is $O(\varepsilon^4)$, strictly
faster than the Cauchy $O(\varepsilon^2)$.
The Cauchy blob is therefore the slowest-converging profile,
and all finite-$\sigma^2$ profiles inherit the stability
conclusions at $O(\varepsilon^2)$ or better.
\end{proof}

The Hessian projection $P_m$ and the total correction
$c_m = -(N^2{-}1)/12 + P_m$ (specific to the Cauchy blob
$h_\varepsilon = -\tfrac12\ln(d^2 + 2\varepsilon^2)$) are exact
rational numbers, computed as follows.  Differentiating the Cauchy
blob pair Hessian with respect to $\varepsilon^2$ at $\varepsilon=0$
gives the per-pair correction:
\begin{equation}\label{eq:blob_pair_correction}
  \delta H_{jj,xx} = \frac{2}{d_{jk}^6}\bigl((\Delta y)^2 - 3(\Delta x)^2\bigr),
  \;\;\;
  \delta H_{jk,xx} = -\delta H_{jj,xx},
\end{equation}
with analogous expressions for $yy$ and $xy$ (see proof of
Proposition~\ref{prop:blob}).  Here $(\Delta x, \Delta y, d_{jk})$
are the pair separation components.  The critical radial Fourier
mode is $v_k \propto z_k\cos(2\pi mk/N)$ with $m = \lfloor N/2\rfloor$.
The projection
\begin{equation}\label{eq:Pm_formula}
  P_m = v_m^\top\,\delta\nabla^2\!H\,v_m
\end{equation}
is a finite quadratic form: a sum over $N(N{-}1)$ pair
contributions, each involving trigonometric values at $N$-th roots
of unity.  Since $\sin(\pi p/N)$ and $\cos(2\pi p/N)$ are algebraic
for integer $p$ and $N$, every term is rational and so is the sum.
Explicit values (verified to machine precision;
see Code Availability):
\begin{center}
\begin{tabular}{lcccccc}
\toprule
$N$ & 3 & 4 & 5 & 6 & 7 & 8 \\
\midrule
$-(N^2{-}1)/12$ & $-\frac23$ & $-\frac54$ & $-2$ &
  $-\frac{35}{12}$ & $-4$ & $-\frac{21}{4}$ \\[3pt]
$P_m$ & $-\frac13$ & $+\frac14$ & $+1$ & $+4$ & $+8$ &
  $+\frac{65}{4}$ \\[3pt]
$c_m$ & $-1$ & $-1$ & $-1$ & $+\frac{13}{12}$ & $+4$ & $+11$ \\
\bottomrule
\end{tabular}
\end{center}

\noindent
The Hessian projection $P_m$ changes sign between $N = 3$
($P_3 = -1/3$) and $N = 4$ ($P_4 = +1/4$): for the triangle, the blob
softening reduces the nearest-neighbor curvature (all pairs are at the
same short distance $d = \sqrt{3}$), while for the square and larger
polygons, the long-range pairs increasingly dominate the projection
and $P_m > 0$.  Despite this sign change,
$c_m = -1$ universally for $N \le 5$ because the Lagrange multiplier
shift $-(N^2{-}1)/12$ dominates $P_m$ in this range.  For $N \ge 6$,
$P_m$ grows faster and overwhelms the shift.

By profile independence (proved above), all
finite-$\sigma^2$ blobs have smaller corrections,
so the stability conclusions hold \emph{a fortiori}.
Three consequences follow:

\emph{(i) Stability preserved.}  For $N \le 5$:
$\lambda_0 \ge 1$ and $c = -1$, so
$\lambda(\varepsilon) = \lambda_0 - \varepsilon^2 > 0$
for $\varepsilon < 1$.
For $N = 6$: $c > 0$, so $\lambda(\varepsilon) > \lambda_0$, the blob \emph{enhances} stability.

\emph{(ii) Marginal mode resolved.}  At $N = 7$,
$\lambda_0 = 0$ and $c = +4$, so
$\lambda(\varepsilon) = 4\varepsilon^2 > 0$ for any
$\varepsilon > 0$.  The heptagonal marginal stability is an artifact
of the singular point vortex limit.

\emph{(iii) Instability partially compensated.}  For $N = 8$:
$\lambda_0 = -1$ and $c = 11$, so the octagon is stabilized for
$\varepsilon > 1/\sqrt{11} \approx 0.30$.

\subsubsection{Full infinite-dimensional convergence}

The full infinite-dimensional stability analysis (including blob shape
deformations and center-shape coupling) is given in
Appendix~\ref{app:convergence}.  The main result
(Proposition~\ref{prop:inertia}) is that for sufficiently concentrated
blobs, the full Hessian inertia equals the point vortex Havelock inertia:
no eigenvalue is lost, no sign is changed, and no new instability arises
from the infinite-dimensional shape degrees of freedom.
Corollaries~\ref{cor:inertia-N6} and~\ref{cor:inertia-N7} give the
explicit concentration bounds for $N \le 6$ and the marginal-mode
resolution for $N = 7$.

\subsection{The Rossby bridge (Saturn regime)}
\label{sec:rossby-bridge}

For continuous jet meanders ($\beta R^3/\kappa \gg 1$), the bridge
is not vortex concentration but \emph{Rossby wave quantization}.
The QGPV linearized about a zonal jet $U(y)$ gives the Rayleigh--Kuo
equation:
\[
  (U - c)\bigl(\hat\psi'' - k^2\hat\psi\bigr)
  + (\beta - U'')\hat\psi = 0.
\]
For a stationary wave ($c = 0$) at the jet center ($y = 0$) of a
Gaussian jet with $U''(0) = -U_{\max}/\sigma^2$:
\[
  \hat\psi'' + \Bigl[\frac{\beta_{\mathrm{eff}}}{U_{\max}}
  - \frac{n^2}{R^2}\Bigr]\hat\psi = 0,
  \qquad
  \beta_{\mathrm{eff}} = \beta + \frac{U_{\max}}{\sigma^2},
\]
where the effective PV gradient $\beta_{\mathrm{eff}}$ includes the
jet curvature (which dominates $\beta$ by a factor $\sim\!13$ for
Saturn).  The stationary azimuthal wavenumber is
\begin{equation}\label{eq:rossby-bridge}
  n^* = R_{\mathrm{S}}\cos\varphi
  \,\sqrt{\beta_{\mathrm{eff}}/U_{\max}},
\end{equation}
where $R_{\mathrm{S}}\cos\varphi$ is the ring radius at latitude
$\varphi$ ($\approx 1.31 \times 10^7\;\mathrm{m}$ for Saturn at
$76^\circ$\,N).  The nearest integer $n = \mathrm{round}(n^*)$ is
the selected polygon wavenumber.  For Saturn:
$\beta_{\mathrm{eff}}/U_{\max} \approx 1.72 \times
10^{-13}\;\mathrm{m}^{-2}$, giving $n^* \approx 5.4$ from the
leading-order dispersion relation.  The observed $n = 6$ is obtained
in numerical eigenvalue computations of the full Rayleigh--Kuo
equation \citep{Morales2011}; the discrepancy between the
leading-order $n^* \approx 5.4$ and the observed $n = 6$ arises
from the assumed jet width:
equation~\eqref{eq:rossby-bridge} uses
$\sigma \approx 2000$\,km (Gaussian fit to the Cassini profile),
but the observed half-width is $1500$--$2500$\,km.
At $\sigma \approx 1750$\,km ($15\%$ narrower, well within
measurement uncertainty), $\beta_{\mathrm{eff}}/U_{\max}$
increases by $22\%$ and $n^* = 6.0$ exactly.
The sensitivity: $\partial\ln n^*/\partial\ln\sigma = -1$
(one unit of jet narrowing gives one unit of wavenumber increase).

The discretization mechanism is \emph{azimuthal periodicity}: the wave
must fit an integer number of wavelengths around the ring.

The connection to the point vortex energy landscape proceeds through
the \emph{critical layer} of the saturated wave.  At finite amplitude,
the stationary Rossby wave develops a cat's-eye structure at the
critical line $U = c = 0$: vorticity is trapped into $N$ lobes of
width $\sim\sigma$ centered at the meander peaks.  As the jet sharpens
($\sigma/R \to 0$ at fixed wave amplitude), these lobes concentrate
into $N$ isolated patches near the $N$-gon vertices.

\begin{remark}[Cat's-eye Hessian convergence]
\label{rem:catseye}
The cat's-eye vorticity decomposes as
$\omega = \omega_{\mathrm{lobe}} + \omega_{\mathrm{braid}} +
\omega_{\mathrm{bg}}$, where the lobes are the trapped fluid inside
the separatrices, the braids are thin filaments along the separatrix
itself, and the background is the passing fluid outside the cat's-eye.

Two conditions are required for Proposition~\ref{prop:inertia} to
apply:
\begin{itemize}
  \item[(i)] \textbf{Lobe concentration.}  The lobes satisfy the
    concentration condition~\eqref{eq:concentration-bound} with
    $\varepsilon \sim \sigma$ and the centroids converge to the
    $N$-gon vertices.  This is guaranteed by $\mathbb{Z}_N$ symmetry.
  \item[(ii)] \textbf{Braid smallness.}  The absolute braid circulation
    $\Gamma_{\mathrm{braid}}$ is negligible compared to the total
    lobe circulation~$\Gamma_{\mathrm{lobe}}$.
\end{itemize}

Both conditions hold for the Stuart vortex
$\Psi = -\sigma^2\ln(\cosh(y/\sigma) - \eta\cos(n\theta))$
($\eta \in (0,1)$; $\eta \to 1$ is point-vortex concentration).

\emph{Proof of lobe dominance.}  The Stuart vortex satisfies
$\omega = C/F^2$ with $F = \cosh(y/\sigma) - \eta\cos(n\theta)$.
Near the elliptic center ($\theta = 0$, $y = 0$):
$F \approx (1{-}\eta) + (y/\sigma)^2/2 + n^2\eta\,\theta^2/2$,
so
\[
  \Gamma_{\mathrm{lobe}} = C\iint_{\mathrm{lobe}} \frac{dA}{F^2}
  \ge C\int_0^\infty \frac{2\pi r\,dr}{[(1{-}\eta)+r^2/2]^2}
  = \frac{2\pi C}{1{-}\eta}.
\]
Outside the separatrix, $F \ge 1{+}\eta \ge 1$, so
$\Gamma_{\mathrm{non\text{-}lobe}} = C\!\iint_{\mathrm{outside}} F^{-2}\,dA
\le C\,A_{\mathrm{period}}$, which is bounded independently of~$\eta$.
Therefore
\begin{equation}\label{eq:lobe-fraction}
  \frac{\Gamma_{\mathrm{non\text{-}lobe}}}{\Gamma_{\mathrm{total}}}
  \;\le\; \frac{A_{\mathrm{period}}\,(1{-}\eta)}{2\pi + A_{\mathrm{period}}\,(1{-}\eta)}
  \;=\; O(1{-}\eta).
\end{equation}
The lobe fraction $\Gamma_{\mathrm{lobe}}/\Gamma_{\mathrm{total}}
= 1 - O(1{-}\eta)$ as $\eta \to 1$.  Numerically: at $\eta = 0.99$,
lobes capture $>99\%$ of the total circulation.

Combining with Proposition~\ref{prop:inertia}:
the cat's-eye satisfies the blob bridge conditions
when $\sigma/d_{\min} < \delta(N)/\!\sqrt{2}$
and $1 - \eta < c\,\delta(N)^2$.
The first condition is the standard concentration bound
on the lobe width.
The second condition holds because the braid perturbation
is $O(1{-}\eta)$ of the lobe system
(equation~\eqref{eq:lobe-fraction}).
Both conditions are satisfied for~$\eta$ sufficiently
close to~$1$ (the point-vortex limit).
Therefore, by Proposition~\ref{prop:inertia},
the constrained Hessian inertia of the cat's-eye
equals the Havelock inertia.
The numerical verification
($\eta \in [0.3, 0.99]$; see Code Availability) confirms this.

The same argument applies to the Prandtl--Batchelor
cat's-eye, the physically relevant model for Saturn's
hexagon, where the lobe dominance is \emph{exact}:
the PB theorem \citep{Batchelor1956} gives
$\omega = \omega_{\mathrm{lobe}}$ (constant) inside each
separatrix and $\omega = \omega_{\mathrm{bg}}$ outside.
The braid (the separatrix curve $\psi = \psi_{\mathrm{sep}}$)
is a single streamline of zero area, so
$\Gamma_{\mathrm{braid}}
= \int_{\mathrm{separatrix}}\!\omega\,dA = 0$ exactly.
Every unit of circulation is in either lobe or background;
lobe dominance~\eqref{eq:lobe-fraction} holds with
$1{-}\eta$ replaced by~$0$.
See Code Availability.
\end{remark}

The Saturn bridge thus arrives at the \emph{same} energy landscape as
the Jupiter bridge, through a different limiting procedure: Saturn
concentrates a continuous jet meander into $N$ patches (jet width
$\sigma \to 0$), while Jupiter concentrates $N$ discrete cyclones
into point vortices (blob radius $\varepsilon \to 0$).
Both limits converge to the Havelock spectrum.

%% ====================================================================
\section{Finite Rossby deformation radius: the $K_0$ interaction}
\label{sec:deformation}

The $K_0$ interaction provides a second example of the RG flow
toward the logarithmic fixed point
(Paper~I, Proposition~2.1): the exponential tail
$K_0(r/R_d) + \ln r = O(e^{-r/R_d})$ is an irrelevant perturbation
(faster than any power law), and the limit $R_d \to \infty$ is the
flow to the fixed point.

In quasi-geostrophic dynamics with finite Rossby deformation radius
$R_d$, the Green's function is $K_0(r/R_d)$ (modified Bessel function),
which recovers $-\ln r$ for $r \ll R_d$.  The interaction $h(r) =
K_0(r/R_d)$ satisfies $h'(r) = -K_1(r/R_d)/R_d < 0$ and
$(rh'(r))' = rK_0(r/R_d)/R_d^2 > 0$ (class~(iii)).  Numerical
computation of the constrained Lagrangian Hessian as a function of
$\eta = R/R_d$ (polygon ring radius over deformation radius) reveals
an $N$-dependent \emph{polygon-suppression transition}:

\begin{center}
\begin{tabular}{lccl}
\toprule
$N$ & Barotropic $\lambda_{\min}$ & Critical $R/R_d$ & Status \\
\midrule
$\le 5$ & $\ge 1$ & none & Stable for all $R/R_d$ \\
$6$ & $1/2$ & $\approx 1.29$ & Suppressed for $R/R_d > 1.29$ \\
$7$ & $0$ & $\approx 0$ & Unstable for any finite $R/R_d$ \\
$\ge 8$ & $< 0$ & N/A & Unstable (barotropic) \\
\bottomrule
\end{tabular}
\end{center}

\noindent
Polygons with $N \le 5$ are robust against stratification effects.
The hexagon ($N = 6$) is the first polygon sensitive to $R_d$: it
requires $R/R_d < 1.29$.  The heptagon ($N = 7$), neutrally stable
in the barotropic limit, is destabilized by any finite
$R/R_d$, providing an additional mechanism for its rarity beyond the
quartic confinement.

For Saturn, the barotropic deformation radius
($\sim\!2000$--$3000\,$km) gives $R/R_d \approx 8$--$12$ at the
hexagonal jet, placing it above the critical ratio.  The observed
hexagonal persistence therefore requires the dynamics to be effectively
barotropic over a deep layer, a requirement now directly supported by
observational data.  Cassini Grand Finale gravity measurements constrain
Saturn's differential rotation to a depth of $8{,}000$--$9{,}000$\,km
\citep{Iess2019}, implying a coherent barotropic column of order
$10^4$\,km, far exceeding the hexagon's vertical scale.  Saturn is
therefore the best-supported case in our sample for the applicability
of the 2D logarithmic Green's function: the columnar depth is
observationally established.

\paragraph{Observable prediction: hexagon suppression signatures.}
The critical ratio $R/R_d \approx 1.29$ for $N = 6$ yields a
testable prediction.  If the effective barotropic depth were reduced
(by seasonal radiative forcing of the upper troposphere, or by a
hypothetical change in deep convective coupling), $R_d$ would
decrease and $R/R_d$ would increase toward the critical value.
Approaching the threshold, the framework predicts three observable
signatures: (1)~the minimum Havelock eigenvalue
$\lambda_{\min}(\eta) \to 0$, reducing the hexagonal mode's
restoring frequency and producing a \emph{longer oscillation period}
for meridional meander perturbations;
(2)~the hexagonal amplitude $\varepsilon_{\mathrm{obs}}$ would
\emph{grow}, since the restoring force weakens while convective
forcing remains constant; and
(3)~at $R/R_d > 1.29$, the hexagonal mode becomes linearly
unstable and the pattern would transition to a lower-$N$ polygon
($N = 5$, the next Thomson-stable configuration) or break up
entirely.
Saturn's hexagon has been observed continuously since Voyager~1
(1980--81); Cassini monitored it through northern spring equinox
(2009) to summer solstice (2017) with no sign of weakening
\citep{Antunano2015, SanchezLavega2014}.  This four-decade
persistence across half a Saturnian year places an empirical lower
bound on the effective barotropic depth: seasonal radiative forcing
does not reduce $R_d$ enough to approach the critical ratio, consistent
with the deep (${\sim}\,10^4\;\mathrm{km}$) barotropic column
implied by Cassini gravity.

\paragraph{Why the leading $K_0$ correction does not shift stability.}
The small-$r$ expansion $K_0(r/R_d) = -\ln(r/R_d) - \gamma_E
+ r^2/(4R_d^2) + O(r^4)$ shows the leading correction to the
logarithm is $g(r) = r^2/(4R_d^2)$, a perturbation of
RG exponent $\alpha = 2$.  By the isotropic-marginality result
(Paper~I, Proposition~\ref{I-prop:alpha-2-marginal}), the pair
Hessian anisotropy $g'' - g'/r = \alpha(\alpha{-}2)\,r^{\alpha-2}$
vanishes at $\alpha = 2$: the Hessian perturbation is
purely isotropic ($\propto I$) and is exactly compensated by
the rotation-frequency shift $\delta\Omega$.
Consequently $\delta\lambda_m = 0$ for all binding modes, the
$r^2/(4R_d^2)$ term shifts the precession rate of the polygon but
does \emph{not} shift its stability eigenvalues.
The first correction to stability from $K_0$ is $O(r^4/R_d^4)$
($\alpha = 4$, which is \emph{not} isotropic-marginal), and this
is numerically small for $R/R_d \ll 1$.
This explains why the barotropic ($R_d \to \infty$) stability
predictions remain accurate even at moderate $R/R_d$: the
leading deformation-radius correction is stability-neutral.

\paragraph{Jupiter polar cyclones.}
The situation is different, and the 2D justification
requires separate examination.  Juno gravity measurements
\citep{Kaspi2018} constrain Jupiter's differential rotation to
$\sim\!3{,}000$\,km depth (${\approx}\,4\%$ of $R_J$), far shallower
than Saturn's.  Jupiter's polar cyclones therefore do \emph{not} rest
on deep barotropic Taylor columns; the relevant 2D dynamics is that of
a shallow weather layer above a near-solid-body rotating deep interior.
Two limiting cases bracket the problem:

\textit{Deep weather-layer limit ($H_{\rm eff} \sim 3{,}000$\,km).}
Taking the full differential-rotation depth as the effective layer
and using the equivalent-barotropic deformation radius
$R_d = \sqrt{g\,H_{\rm eff}}/f \approx \sqrt{25 \times 3\times10^6}\,/\,
(2\Omega_J\sin 84^\circ) \approx 25{,}000$\,km, we obtain
$R/R_d \approx 0.3 \ll 1$.  In this limit
$K_0(r/R_d) \to -\ln r + \mathrm{const}$, so the logarithmic
Green's function applies, though the physical justification is
a shallow equivalent-barotropic layer rather than a deep Taylor column.

\textit{Tropospheric limit ($H_{\rm eff} \sim 30$--$100$\,km).}
If the active polar-cyclone layer is shallower, reduced gravity
$g' \sim 0.1$--$1$\,m\,s$^{-2}$ gives
$R_d \sim 2{,}000$--$6{,}000$\,km and $R/R_d \sim 1$--$4$.
The $N = 8$ ring lies in the finite-$R_d$ regime.  Computing the
$\kappa_0$-shifted constrained Hessian eigenvalue for the $K_0$
interaction with $\kappa_0/\kappa = 0.7$ (the Juno central estimate):
\begin{center}
\begin{tabular}{cccc}
\toprule
$N$ & $\kappa_0/\kappa$ & $(R/R_d)_{\mathrm{crit}}$ & Tropospheric stable? \\
\midrule
$6$ & $0.7$ & $17.2$ & yes ($\gg 4$) \\
$7$ & $0.7$ & $4.2$  & marginal \\
$8$ & $0.7$ & $\mathbf{0.79}$ & \textbf{no} ($< 1$) \\
$9$ & $0.7$ & $0.10$ & no \\
\bottomrule
\end{tabular}
\end{center}
Here $(R/R_d)_{\mathrm{crit}}$ is the ratio at which the minimum
constrained eigenvalue crosses zero: the ring is stable for
$R/R_d < (R/R_d)_{\mathrm{crit}}$ and unstable above.  For $N = 8$
the critical ratio is $(R/R_d)_{\mathrm{crit}} \approx 0.79$, which
lies \emph{below} the tropospheric range $R/R_d \in [1, 4]$.
\textbf{The $N = 8$ prediction therefore requires the deep
weather-layer limit} ($R/R_d \lesssim 0.8$, i.e.\
$R_d \gtrsim 1.3\,R \approx 9{,}700$\,km); it does not survive
into the tropospheric regime.

The effective layer depth is not directly constrained by current Juno
data, making Jupiter the \emph{least} observationally supported case
for the 2D assumption, in contrast to Saturn.  A direct measurement
of the depth of the polar cyclone roots, analogous to the gravity
constraints on the jets, would resolve this ambiguity.  The
computation above sets a quantitative target: $R_d > 9{,}700$\,km
is necessary for $N = 8$ stability under the $K_0$ interaction.

The critical ratio is defined implicitly by the equation
$\bar\lambda(\eta_{\mathrm{crit}}, N) = 0$, where the dimensionless
eigenvalue $\bar\lambda(\eta, N) = R^2 \cdot \lambda_{\min}(\eta)$
interpolates smoothly from the Havelock value
$\bar\lambda(0, N) = (N{-}1) - \lfloor N/2\rfloor
(N{-}\lfloor N/2\rfloor)/2$ to zero.

The Lagrange multiplier is determined by the stationarity condition
$\partial H/\partial R + \mu_L\,\partial L/\partial R = 0$.  For
$H = -\sum_{j<k}K_0(d_{jk}/R_d)$ with $d_{jk} = 2R\sin(\pi p/N)$
and $L = NR^2$, the radial derivative of each pair contributes
$-K_0'(d_{jk}/R_d)\cdot 2\sin(\pi p/N)/R_d
= K_1(d_{jk}/R_d)\cdot 2\sin(\pi p/N)/R_d$
(using $K_0' = -K_1$).  Projecting onto the constraint gradient
$\partial L/\partial R = 2NR$ and solving gives
\begin{equation}\label{eq:mu-k0}
  -2\mu_L = \frac{\eta}{R}\sum_{p=1}^{N-1}
  \frac{(1-\cos\theta_p)\,K_1(2\eta\sin(\pi p/N))}
       {2\sin(\pi p/N)},
  \quad \theta_p = 2\pi p/N,
\end{equation}
where $\eta = R/R_d$ and the factor $(1{-}\cos\theta_p)$ arises from
the Fourier decomposition of the pairwise Hessian along the radial
mode.  As $\eta \to 0$: $K_1(x) \sim 1/x$ for small $x$, so each
summand reduces to
$(1{-}\cos\theta_p)/(4\eta\sin^2(\pi p/N)) \cdot (\eta/R)$
$= (1{-}\cos\theta_p)/(4R\sin^2(\pi p/N))$.  Summing over $p$ using
the Havelock identity $\sum_p (1{-}\cos\theta_p)/(4\sin^2(\pi p/N))
= (N{-}1)/2$ recovers the flat-plane value $-2\mu_L = (N{-}1)/(2R^2)$.  For $N = 6$, the critical mode $m = 3$
has zero Havelock weight at $p = 2, 4$ (since
$\cos(2\pi \cdot 2 \cdot 3/6) = 1$), so the critical eigenvalue
depends on $K_0$ and $K_1$ at only \emph{two} distinct arguments,
$\eta$ and $2\eta$, rather than three.

No closed-form analog of the Havelock identity is known for $K_0$.
The orthogonality mechanism that collapses the logarithmic sum to
$m(N{-}m)/2$ relies on the identity $\sum_p e^{2\pi i mp/N}\ln
|2\sin(\pi p/N)| = 0$, which has no counterpart for $K_0$:
the Bessel values $K_0(2\eta\sin(\pi p/N))$ are transcendental at
nonzero algebraic arguments by Siegel's theorem on $E$-functions,
but this does not preclude the full weighted sum from having a
closed form.  Whether $\sum_p (\cos\theta_{pm})
K_0(2\eta\sin(\pi p/N))$ admits a closed form for any $N$ and $m$
is an open question; no such form has been found, and all stability
computations for $K_0$ interactions are carried out numerically.
A perturbative expansion
$K_0(x) \approx -\ln x + O(x^2)$ gives $\eta_{\mathrm{crit}} \approx
0.9$ for $N = 6$ (30\% below the exact value) because the critical
ratio lies in the non-perturbative regime where $K_0$ departs
significantly from $-\ln$.  The transition occurs precisely where the
exponential decay of $K_0$ (destroying long-range correlations between
non-adjacent vortices) overtakes the logarithmic regime, making the
suppression inherently non-perturbative.  The implicit equation is
nevertheless well-posed, and the critical $\eta$ is uniquely determined
by bisection for each $N$.

\paragraph{The two phase boundaries.}
For $K_0$, two distinct boundaries exist in the $(N, \eta)$ plane:
(1)~the \emph{spiral sign change} $\bar H''(0) = 0$, where the
positive $A_m$ term from $(rh')' > 0$ overwhelms the negative $B_m$
term; and (2)~the \emph{constrained stability loss}
$\bar\lambda_{\min} = 0$, where the Lagrange multiplier shift
consumes the stability margin.  The constrained boundary
($\eta \approx 1.29$ for $N = 6$) is always \emph{inside} the
spiral boundary ($\eta \approx 3.2$ for $N = 6$) because the
Lagrange shift reduces the margin before $H''$ changes sign.
Physically, the polygon requires long-range correlations between
non-adjacent vortices.  When $\eta \gg 1$ and the $K_0$ interaction
decays exponentially, only nearest-neighbor correlations survive and
the polygon loses its energetic advantage.

For comparison, the Cauchy-blob regularization
$-\tfrac12\ln(d^2 + 2\varepsilon^2)$ has $H'' < 0$ for \emph{all}
$\varepsilon$, the soft-core never flips the spiral sign.  The
Gaussian-screened interaction $-\ln(d)\cdot e^{-d^2/L^2}$
shows an earlier transition ($\eta \approx 0.9$ for $N = 6$) due to
its faster-than-exponential decay.

\paragraph{The constrained stability boundary.}
The interaction class determines $N_{\mathrm{crit}}$, defined here as
the largest $N$ for which the \emph{constrained} Lagrangian Hessian
$\bar\lambda_{\min} \ge 0$ (i.e.\ the polygon is a constrained energy
minimum).  This is distinct from the spiral-direction boundary
$H''(0) = 0$, which occurs at larger $\eta$ (e.g.\ $\eta \approx 3.2$
vs.\ $1.29$ for $N = 6$ with $K_0$; see the two phase boundaries above):
\begin{center}
\begin{tabular}{lll}
\toprule
Class & Condition & $N_{\mathrm{crit}}$ (constrained) \\
\midrule
(i) & $(rh')' < 0$ & $\ge 7$ (increases with $|(rh')'|$) \\
(ii) & $(rh')' = 0$ (only $-\ln r$) & $= 7$ exactly \\
(iii) & $(rh')' > 0$ (e.g.\ $K_0$) & $\le 7$
  (decreases with screening) \\
\bottomrule
\end{tabular}
\end{center}
\noindent
Numerical computation confirms: for $h = -r^\alpha$ (class~(i)),
$N_{\mathrm{crit}} = 7$ at $\alpha \to 0$ and increases with $\alpha$
($N_{\mathrm{crit}} = 8$ at $\alpha = 1$, $N_{\mathrm{crit}} = 14$
at $\alpha = 2$).  At $N = 7$, $\lambda_{\min}(\alpha) > 0$ for all
$\alpha > 0$ tested, with
$\dd^2\lambda_{\min}/\dd\alpha^2 = +2.5 > 0$ at $\alpha \to 0$:
the marginal stability at $N = 7$ is \emph{unique to the logarithmic
interaction}, and any class~(i) interaction makes the heptagon strictly
stable.  For $K_0(r/R_d)$ (class~(iii)), the dimensionless eigenvalue
$\bar\lambda(\eta)$ decreases monotonically in $\eta$ for each Fourier
mode (verified for $N = 5, 6$), consistent with the physical expectation
that screening weakens the stabilizing long-range correlations.
For the $K_0$ interaction,
$N_{\mathrm{crit}} = 6$ at $R/R_d = 1$ and decreases to~$5$ at
$R/R_d = 3$.  The logarithmic interaction is the \emph{watershed}:
the critical case between interactions that stabilize polygons
(class~(i)) and interactions that suppress them (class~(iii)).

\begin{proposition}[Min/max duality]\label{prop:duality}
For any pairwise interaction $h(r)$, the constrained energy minimum
of $H = \sum h(d_{jk})$ and the constrained energy maximum of $-H$
have the same stability boundary $N_{\mathrm{crit}}$.
\end{proposition}

\begin{proof}
The Lagrange multiplier for $-H$ is $\mu_L^{-H} = -\mu_L^H$
(from $-\nabla H = (-\mu_L)\nabla L$).  The Lagrangian Hessian
$\mathcal{L}_{-H} = -\nabla^2\!H + 2\mu_L I = -\mathcal{L}_H$.
The eigenvalues are negated; constrained minimum (all $\ge 0$) of $H$
is equivalent to constrained maximum (all $\le 0$) of $-H$.
\end{proof}

For gravity ($h = -1/r$, $h' > 0$), the $N$-gon is a constrained
energy maximum.  The duality gives
$N_{\mathrm{crit}}^{\mathrm{gravity}} =
N_{\mathrm{crit}}^{\mathrm{Coulomb}}$.  Numerical computation of
the Coulomb ($h = +1/r$) constrained Hessian gives
$N_{\mathrm{crit}} = 5$: the minimum eigenvalue is $+0.73$ at $N = 5$
and $-0.02$ at $N = 6$.  This is consistent with Maxwell's (1859)
result that gravitating rings require a central mass for
$N \ge 7$; the one-unit difference ($N_{\mathrm{crit}} = 5$ vs
Maxwell's $6$) arises because our analysis does not include the
central mass (Maxwell's stability includes the stabilizing effect
of the central body).

Maxwell (1859) obtained $N_{\mathrm{crit}} = 6$ for the Coulomb ring \emph{with} a central mass; without the central mass, $N_{\mathrm{crit}} = 5$. The central mass adds a uniform eigenvalue shift $+\kappa_0/R^2$ that stabilizes $N = 6$.

The Coulomb/gravity $N_{\mathrm{crit}} = 5$ is less than the
logarithmic $N_{\mathrm{crit}} = 7$, confirming the class~(iii)
reduction: the $1/r$ interaction has $(rh')' = 1/r^2 > 0$, placing
it in the destabilizing regime relative to $-\ln r$.


%% ====================================================================
\section{Oblateness and the ice giants}
\label{sec:oblateness}
%% ====================================================================

\subsection{Saturn's oblateness determines hexagon existence}
\label{sec:saturn-oblateness}

\begin{figure}[t]
\centering
\includegraphics[width=0.75\textwidth]{../figures/fig_oblateness_signflip.pdf}
\caption{The oblateness sign-flip: Saturn's oblateness
($f = 0.098$) reverses the $N = 6$ eigenvalue from
$\lambda_3 = -0.12$ (unstable) to $+0.04$ (stable),
determining whether the hexagon exists.}
\label{fig:oblateness}
\end{figure}

The preceding sections analyzed stability on
constant-curvature surfaces and with finite deformation
radius.  A third physical correction, planetary
oblateness, has a qualitatively different effect: it can
reverse the sign of an eigenvalue.

The Rossby analysis of \S\ref{sec:rossby-bridge} treats
Saturn as a sphere.
Saturn is an oblate spheroid with flattening
$f = (a{-}b)/a = 0.098$
($a = 60{,}268$\,km equatorial, $b = 54{,}364$\,km polar).
The oblateness changes the \emph{sign} of the Thomson stability
eigenvalue for $N = 6$: it is not a small correction but the
mechanism that determines whether the hexagon can exist.

The Gaussian curvature at geodetic latitude~$\varphi$ is
$K(\varphi) = 1/(M(\varphi)\,N(\varphi))$ where
$M = a(1{-}e^2)/(1{-}e^2\sin^2\!\varphi)^{3/2}$ and
$N = a/(1{-}e^2\sin^2\!\varphi)^{1/2}$.
On an oblate body, $K$ is largest at the equator and
smallest at the pole, the pole is flatter.
For Saturn: $K_{\mathrm{eq}}/K_{\mathrm{pole}} = a^4/b^4 = 1.51$.
Since positive curvature destabilizes
(Paper~I, curvature dichotomy), the pole is
the most stable point on the spheroid.

Substituting the local curvature into the spherical
$C_1$ formula:
\begin{center}
\begin{tabular}{lcc}
\toprule
Model & $\lambda_3(N{=}6)$ & Stable? \\
\midrule
Mean sphere ($R_{\mathrm{mean}} = 58{,}300$\,km) &
  $-0.12$ & No \\
Oblate spheroid at $78^\circ$N &
  $+0.04$ & Yes \\
\bottomrule
\end{tabular}
\end{center}
The oblateness flips $\lambda_3$ from $-0.12$ to $+0.04$.
On a mean sphere, the hexagon is Thomson-unstable and
should not exist; on the oblate spheroid, it is
marginally stable.

\begin{proposition}[Oblateness-stabilized hexagon]
\label{prop:oblateness}
On an oblate spheroid with flattening~$f$, the binding
eigenvalue $\lambda_3(N{=}6)$ at the pole exceeds its mean-sphere
value by
\begin{equation}\label{eq:oblateness-shift}
  \delta\lambda_3 \approx (N{-}1)\cdot
  \frac{6 f\,\xi}{(1{+}\xi)^2}
\end{equation}
at leading order in~$f$, where $\xi = K_{\mathrm{mean}} R^2$.
For Saturn ($f = 0.098$, $\xi = K_{\mathrm{mean}} R^2
\approx (15{,}000/58{,}300)^2 = 0.066$):
the leading-order formula gives $\delta\lambda_3 \approx 0.17$;
the exact finite-difference computation
$C_1(\xi{+}\delta\xi){-}C_1(\xi)$ with
$\delta\xi = \xi[(1{-}f)^3{-}1]$
gives $\delta\lambda_3 = 0.16$,
the residual being $O(f^2)$ corrections.
This reverses the sign from $-0.12$ to $+0.04$
(comfortably positive, not a marginal flip).
\end{proposition}

\begin{proof}[Proof sketch]
The Gaussian curvature at the pole of an oblate spheroid
is $K_{\mathrm{pole}} = b^2/a^4
= K_{\mathrm{mean}}\,(1{-}f)^3$
(since $K_{\mathrm{mean}} \approx 1/(ab)$
and $b/a = 1{-}f$).
The sphere threshold formula
$C_1^{S^2}(\xi) = (N{-}1)(1{-}\xi)/(1{+}\xi)$
evaluated at $\xi_{\mathrm{pole}} = K_{\mathrm{pole}} R^2$
gives, relative to $\xi_{\mathrm{mean}}$:
$\delta C_1 = -(N{-}1) \times 2\,\delta\xi/(1{+}\xi)^2
+ O(\delta\xi^2)$,
where $\delta\xi = \xi_{\mathrm{pole}} - \xi_{\mathrm{mean}}
= \xi_{\mathrm{mean}}[(1{-}f)^3 - 1]
\approx -3f\,\xi_{\mathrm{mean}}$.
Since the pole is flatter ($K_{\mathrm{pole}} < K_{\mathrm{eq}}$,
hence $\delta\xi < 0$),
$\delta C_1 > 0$: the eigenvalue increases.
The leading-order approximation gives
$\delta\lambda_3 \approx (N{-}1) \times 6f\xi/(1{+}\xi)^2$,
equation~\eqref{eq:oblateness-shift}.
The exact expansion retains the $f^2$ and $f^3$ terms
from $(1{-}f)^3 - 1 = -f(3 - 3f + f^2)$,
giving
$\delta\lambda_3 = (N{-}1) \cdot 2\xi f(3 - 3f + f^2)/(1{+}\xi)^2$.
\end{proof}

\noindent
The critical latitude $\varphi_{\mathrm{crit}}$ where $\lambda_3 = 0$
is $\approx 72^\circ$N for Saturn.
The hexagon at $78^\circ$N is $6^\circ$ above this boundary, comfortably
in the stable region.

\begin{remark}[Crystallographic uniqueness of $N = 6$]
\label{rem:crystallographic-N6}
Three equivalent conditions single out $N \in \{1, 2, 3, 4, 6\}$
among all positive integers:
(i)~$\varphi(N) \le 2$, i.e.\ $[\mathbb{Q}(\zeta_N):\mathbb{Q}] \le 2$;
(ii)~$2\cos(2\pi/N) \in \mathbb{Z}$, the crystallographic
restriction for rotations preserving a $2$D lattice; and
(iii)~$\mathrm{Li}_s(\zeta_N)$ admits a closed-form decomposition
over $\mathbb{Q}$ in terms of $\zeta(s)$ and an $L$-function with
rational Euler factors at primes dividing~$N$.
Among these, $N = 6$ is the unique composite case, with
Euler factors at both primes dividing $6 = 2 \cdot 3$.
Non-crystallographic $N$ (including $N = 5, 7, 8, \ldots$)
can still support finite rings but require incommensurate
arithmetic (golden ratio for~$N = 5$, Pell units for~$N = 7$):
see \citealt*{Shechtman1984} for quasi-crystal realizations.
On planetary scales, the stable polygons realized in
observation ($N = 6$ on Saturn, $N = 8$ with central vortex
on Jupiter) are precisely the smallest values that combine
crystallographic compatibility ($N = 6$) or stability via
central-vortex augmentation ($N = 8$); larger stable
non-crystallographic polygons are geometrically possible
but not observed at planetary scales.
The hexagonal lattice energy is given in closed form by the
Epstein zeta factorization
$\zeta_{A_2}(s) = \zeta(s)\,L(s, \chi_{-3})$
(Hecke; the $A_2$ lattice is the Eisenstein integers
$\mathbb{Z}[\omega]$ with $\omega = e^{2\pi i/3}$).
\end{remark}

At non-polar latitudes, the curvature gradient
$\delta K/K \approx 4.2\%$ across the hexagon ring
couples Fourier modes $m$ and $m \pm 2$.
The coupling amplitude ($\sim 0.19$) exceeds the stability
margin ($|\lambda_3| = 0.04$) by a factor of~$5$.
By the topological protection of the Morse index
(Paper~I), this mode coupling cannot change the net spectral
flow, the hexagon remains stable in the topological sense.

\subsection{Neptune predictions: packing-dominated}
\label{sec:neptune}

Neptune ($R = 24{,}600$\,km, $\Omega = 1.08\times10^{-4}$\,rad\,s$^{-1}$,
$f = 0.017$) is in the vortex-dominated regime:
at $84^\circ$N,
$\beta_N = 2\Omega_N\cos(84^\circ)/R_N
\approx 9.2 \times 10^{-13}$\,m$^{-1}$s$^{-1}$
and $R_{\mathrm{ring}} \approx 2{,}570$\,km, giving
$\beta R^3 \approx 1.6 \times 10^{7}$\,m$^2$\,s$^{-1}$.
For Jovian-scaled cyclones
($V_{\max} \approx 80$--$100$\,m\,s$^{-1}$,
$r_c \approx 240$\,km),
$\kappa \approx 1.2$--$1.5 \times 10^{8}$\,m$^2$\,s$^{-1}$,
so $\beta R^3/\kappa \approx 0.10$--$0.13 \ll 1$.
The Rossby stationarity mechanism does not apply;
the relevant bridge is Thomson/Onsager.

The binding constraint is \emph{packing}: Neptune's polar
ring radius ($R_{\mathrm{ring}} \approx 5{,}000$\,km,
estimated from the polar vortex extent) is too small
for many cyclones of Jupiter-like size to fit.
With a $\beta$-drift exclusion factor ${\approx}\,1.3$
(estimated from the spacing analysis of \citealt{Gavriel2021}),
the effective exclusion radius is
$r_{\mathrm{excl}} = 1.3\,R_d$, and the prediction depends
on the deformation radius:
\begin{center}
\small
\begin{tabular}{lccc}
\toprule
$R_d$ (km) & $r_{\mathrm{excl}}/R_{\mathrm{ring}}$ &
  $N_{\mathrm{pack}}$ & $N_{\mathrm{pred}}$ \\
\midrule
$1{,}500$ & $0.39$ & $7$ & $\le 7$ \\
$2{,}000$ (moderate) & $0.52$ & $5$ & $5$ \\
$2{,}500$ (Jupiter-like) & $0.65$ & $4$ & $4$ \\
\bottomrule
\end{tabular}
\end{center}
Neptune's strong internal heat flux
($F_{\mathrm{int}} = 0.43$\,W\,m$^{-2}$,
$2.6\times$ absorbed solar) favours vigorous convection,
suggesting $R_d \gtrsim 1{,}000$\,km.
In this regime, the prediction narrows to
$N \in \{3, 4, 5\}$, a triangle, square, or pentagon.

Neptune's oblateness ($f = 0.017$) is negligible:
$\delta\lambda_3 \lesssim 0.01$, far below the
Thomson eigenvalue margin.

This prediction is testable by a Neptune polar orbiter
(Trident-class mission): imaging the polar region at
$\sim 100$\,km resolution would resolve individual
cyclones and determine~$N$ directly.

\emph{Falsifiability}: a Neptune observation of $N \ge 6$
would violate the packing bound for any
$R_d > 1{,}500$\,km.

\subsection{Uranus: no polygon predicted}
\label{sec:uranus}

Uranus (axial tilt $98^\circ$) presents a qualitatively
different case.
The extreme tilt disrupts the Onsager condensation
mechanism by preventing a persistent polar jet:
solar forcing alternates between $42$-year polar day and
$42$-year polar night, and the pole receives
\emph{more} cumulative insolation than the equator
over an orbit.
The prediction is $N = 0$, no polar polygon.
This argument is qualitative: there is no sharp
tilt threshold in the framework at which
the Onsager mechanism switches off.
The $98^\circ$ tilt is invoked as an extreme case
where the persistent differential heating needed for
a polar jet is absent; a quantitative criterion
(e.g.\ minimum jet persistence time as a function
of obliquity) remains to be derived.

Current observations are consistent:
\citet{Akins2023} detect a single polar cyclone at
Uranus's north pole ($N = 1$, below the polygon threshold).
No ring structure has been observed.

Uranus's oblateness ($f = 0.023$) is negligible for
the same reason as Neptune.

\emph{Falsifiability}: a Uranus polar polygon would
challenge the tilt-disruption hypothesis.

\subsection{Why the hexagon is uniquely two-dimensional}
\label{sec:uniquely-2d}

The hexagon ($N = 6$) is marginally stable in $d = 2$
($\lambda_{\min} = +0.5$) but unstable for $d \ge 3$
($\lambda_{\min} = -0.018$ at $d = 3$):
\begin{center}
\small
\begin{tabular}{lcccl}
\toprule
$d$ & $\Delta = (d{-}2)/2$ & $N_{\mathrm{crit}}$ &
  $\lambda_{\min}(N{=}6)$ & Hexagon \\
\midrule
$2$ & $0$ & $7$ & $+0.50$ & Stable \\
$3$ & $1/2$ & $5$ & $-0.02$ & Unstable \\
$4$ & $1$ & $5$ & $-1.08$ & Unstable \\
\bottomrule
\end{tabular}
\end{center}
\noindent
(The $d \ge 3$ eigenvalues use the $\csc^d$ interaction
kernel normalized to match the $d = 2$ self-energy
at the polygon radius; see Code Availability.)

Saturn's hexagon requires the logarithmic interaction
($d = 2$, $\Delta = 0$), the Green's function of
the 2D Laplacian.
The logarithmic interaction survives in a quasi-2D
atmosphere because the weather layer is thin
($H_{\mathrm{eff}} \ll R_{\mathrm{planet}}$) and the
barotropic mode dominates.
This explains why hexagonal features are observed only on
gas giants with thin weather layers and not on
terrestrial planets or in 3D laboratory flows.


%% ====================================================================
\section{The constraint-intersection principle}
\label{sec:constraint-intersection}

\begin{figure}[t]
\centering
\includegraphics[width=0.75\textwidth]{../figures/fig_constraint_intersection.pdf}
\caption{The three constraints (Rossby, Thomson, packing)
intersect to select $N$ on each planet.
Saturn ($N = 6$): Rossby-limited.
Jupiter north ($N = 8$): Thomson-limited.
Jupiter south ($N = 5$): packing-limited.}
\label{fig:constraint-intersection}
\end{figure}

The observed polygon wavenumber $N$ on each planet is determined by the
intersection of three constraints, all descending from the 2D Laplacian.

\begin{proposition}[Constraint-intersection selection]
\label{prop:complementary}
The selected polygon wavenumber is
\begin{equation}\label{eq:constraint-intersection}
  N_{\mathrm{selected}} = \max\bigl\{N :
    N \le N_{\mathrm{Thomson}},\;
    \sin(\pi/N) \ge r_{\mathrm{excl}}/R,\;
    n\;\text{Rossby-stationary}\bigr\},
\end{equation}
where inapplicable constraints are dropped.  The three constraints are:
\begin{enumerate}
  \item \textbf{Rossby stationarity} (derived; for jet meanders):
    the QGPV eigenvalue problem
    $(\nabla^2 + \beta\,\partial_y)\psi = 0$ selects the stationary
    wavenumber $n^*$.
  \item \textbf{Thomson stability} (derived; for discrete cyclone rings):
    $\kappa_0/\kappa \ge \kappa_{\mathrm{crit}}(N)$ gives an upper
    bound $N \le N_{\max}(\kappa_0/\kappa)$; the energy monotonicity
    $H_{N+1} > H_N$ (Proposition~\ref{prop:onsager-n-select}) then
    selects $N = N_{\max}$.
  \item \textbf{Geometric packing} (derived skeleton + empirical factor):
    domain admissibility of $G(\nabla^2)$ on the punctured domain
    gives $\sin(\pi/N) \ge \varepsilon/R$ where $\varepsilon$ is
    the vortex core radius (derived).  The effective exclusion radius
    $r_{\mathrm{excl}} \approx 1.3\,\varepsilon$ includes a
    $\beta$-drift buffer zone from nonlinear vortex dynamics
    \citep{Gavriel2021}; the factor $1.3$ is observational, not
    derivable from the linear operator.
\end{enumerate}
\end{proposition}

Each planetary system activates a different subset of the three
constraints:
\begin{center}
\begin{tabular}{lcccl}
\toprule
System & $N_{\mathrm{Thomson}}$ & $N_{\mathrm{Rossby}}$ & $N_{\mathrm{pack}}$ & Binding \\
\midrule
Saturn (76${}^\circ$N)  & $7$ & $6$ &,  & Rossby \\
Jupiter north           & $8$ &, & $>8$ & Thomson \\
Jupiter south           & $8$ &, & $5$  & Packing \\
\bottomrule
\end{tabular}
\end{center}

\noindent
The two Laplacian-derived facets and their observational supplement are:
\emph{(i)}~spectral positivity of $G(\nabla^2)$: Havelock
eigenvalue $\lambda_m > 0$ $\to$ Thomson bound (derived);
\emph{(ii)}~domain admissibility of $G_\varepsilon(\nabla^2)$:
disc non-overlap $\sin(\pi/N) \ge \varepsilon/R$ $\to$ packing
bound with $r_{\mathrm{excl}} = \varepsilon$ (derived),
the enlargement $r_{\mathrm{excl}} \approx 1.3\,\varepsilon$
at Jupiter south is observational;
\emph{(iii)}~Rossby stationarity of $(\nabla^2 + \beta\,\partial_y)\psi = 0$:
stationary mode $n^*$ $\to$ wave quantization (derived).
The finite-core regularization $G_\varepsilon$ does \emph{not} create
a blob-stable window for $N = 8$: the corrected constrained Hessian
remains negative throughout the admissible domain
$\varepsilon/R < \sin(\pi/8)$.  The $N = 8$ selection therefore
requires the central vortex ($\kappa_0/\kappa \ge \tfrac{1}{2}$);
blob regularization alone does not provide an alternative stabilization
mechanism.

\subsection{Energy monotonicity and maximum-$N$ selection}
\label{sec:energy-monotonicity}

\begin{proposition}[Energy monotonicity and maximum-$N$ selection]
\label{prop:onsager-n-select}
Let $\Gamma = N\kappa$ be the total ring circulation, $L = N\kappa R^2$
the angular impulse (so $R$ is fixed when $\Gamma$ and $L$ are fixed),
and $r = \kappa_0/\kappa$ held fixed as $N$ varies (both the ring and
central circulations scale proportionally).  Denote by
$r_{\mathrm{core}}$ the vortex core radius, the minimum separation at
which the point-vortex model is valid; equivalently, $r_{\mathrm{core}}$
is the unit of length in which the logarithm is taken, so that
$\ln R > 0$ iff $R > r_{\mathrm{core}}$.  Then
\begin{equation}\label{eq:H-monotone}
  H_{N+1} - H_N = \Gamma^2\!\left[\frac{2r-1}{2N(N+1)}\ln R
    + \frac12\!\left(\frac{\ln N}{N} - \frac{\ln(N+1)}{N+1}\right)\right].
\end{equation}
For $r \ge \tfrac{1}{2}$ and $R > r_{\mathrm{core}}$, both terms
are non-negative (first: $2r - 1 \ge 0$ and $\ln R > 0$; second:
$\ln N/N$ is strictly decreasing for $N > e \approx 2.7$), so
$H_{N+1} > H_N$ for all $N \ge 3$.  Therefore the $N$-gon ring
energy is strictly monotone in $N$: the maximum-$H$ configuration
within the Thomson-stable set is the ring with the largest $N$:
\[
  N_{\max}(\kappa_0/\kappa)
    \equiv \max\bigl\{N' : \kappa_{\mathrm{crit}}(N') \le \kappa_0/\kappa\bigr\}.
\]
Any dynamics that preferentially populates higher-energy stable
configurations, whether Onsager negative-temperature equilibrium
($\beta < 0$, where the Boltzmann weight $e^{-\beta H}$ increases
with $H$), sequential vortex injection, or merger of smaller vortex
clusters, selects $N = N_{\max}$.  The stability threshold
$\kappa_{\mathrm{crit}}(8) = \tfrac{1}{2}$ is exactly the value at
which $H_{N+1} > H_N$ becomes unconditional ($R$-independent): at
$r = \tfrac{1}{2}$, term~1 vanishes identically and term~2 remains
strictly positive.
\end{proposition}

\begin{proof}
Use the Hamiltonian
$H = -\kappa^2\!\sum_{j<k}\ln|z_j - z_k| - \kappa_0\kappa\!\sum_j\ln|z_j|$
with $\kappa = \Gamma/N$ and $\kappa_0 = r\kappa = r\Gamma/N$.

\emph{Step 1: Product formula.}
For the regular $N$-gon $z_j = Re^{2\pi ij/N}$, the chord length
between vortices $j$ and $k$ is $|z_j - z_k| = 2R\sin(\pi|j-k|/N)$.
By $\mathbb{Z}_N$ symmetry the pair sum equals $N/2$ times the sum from
one vertex:
\[
  \sum_{j<k}\ln|z_j - z_k|
  = \frac{N}{2}\sum_{s=1}^{N-1}\ln\!\bigl(2R\sin(\pi s/N)\bigr)
  = \frac{N}{2}\bigl[\ln N + (N-1)\ln R\bigr].
\]
Here we used $\sum_{s=1}^{N-1}\ln(2\sin(\pi s/N)) = \ln N$, which
follows from evaluating the polynomial identity $(z^N-1)/(z-1) =
\prod_{k=1}^{N-1}(z - e^{2\pi ik/N})$ at $z = 1$: the product of
moduli is $\prod_{s=1}^{N-1}|1 - e^{2\pi is/N}|
= \prod_{s=1}^{N-1}2\sin(\pi s/N) = N$.

\emph{Step 2: Central vortex term.}
Since $|z_j| = R$ for all $j$: $\sum_j\ln|z_j| = N\ln R$.

\emph{Step 3: Closed form for $H_N$.}
Substituting Steps~1--2 with $\kappa = \Gamma/N$, $\kappa_0 = r\Gamma/N$:
\begin{equation}\label{eq:HN-explicit}
  H_N = -\frac{\Gamma^2}{2N}\ln N
        - \Gamma^2\!\left(\frac{N-1}{2N} + \frac{r}{N}\right)\ln R.
\end{equation}

\emph{Step 4: Difference $H_{N+1} - H_N$.}
Writing $H_{N+1}$ from \eqref{eq:HN-explicit} with $N\to N+1$ and
subtracting:
\begin{align*}
  H_{N+1} - H_N
  &= -\frac{\Gamma^2}{2}\!\left[\frac{\ln(N{+}1)}{N{+}1}
     - \frac{\ln N}{N}\right]
     - \Gamma^2\ln R\!\left[\frac{N/2+r}{N{+}1}
     - \frac{(N{-}1)/2+r}{N}\right].
\end{align*}
The coefficient of $\ln R$ simplifies by partial fractions:
\[
  \frac{N/2+r}{N{+}1} - \frac{(N{-}1)/2+r}{N}
  = \frac{N^2 - (N{-}1)(N{+}1)}{2N(N{+}1)} - \frac{r}{N(N{+}1)}
  = \frac{1-2r}{2N(N{+}1)},
\]
where we used $N^2-(N-1)(N+1) = 1$.  Negating and rearranging
gives~\eqref{eq:H-monotone}.

\emph{Step 5: Sign of each term.}
\emph{Term~2}: $f(x) = \ln x/x$ satisfies $f'(x) = (1-\ln x)/x^2 < 0$
for $x > e$, so $\ln N/N > \ln(N+1)/(N+1)$ for all $N \ge 3$.

\emph{Term~1}: $2r-1 \ge 0$ for $r \ge \tfrac{1}{2}$; and
$\ln R > 0$ for $R > r_{\mathrm{core}}$.  Hence both terms are
non-negative, and term~2 is strictly positive, giving
$H_{N+1} > H_N$.
\end{proof}

\begin{remark}[Physical status of the fixed-$r$ assumption]
\label{rem:fixed-r}
The assumption $r = \kappa_0/\kappa$ fixed as $N$ varies means that
as more cyclones appear in the ring ($N$ increases), each individual
ring vortex becomes weaker ($\kappa = \Gamma/N$) \emph{and the central
vortex weakens proportionally} ($\kappa_0 = r\Gamma/N$).  The
assumption is justified when all vortices are drawn from the same
turbulent condensation process, so their strengths are set by a common
large-scale energy budget.  In the Onsager picture, the
negative-temperature state is reached by the inverse cascade from the
same vorticity reservoir, giving a fixed partition of total circulation
between center and ring.  If instead $\kappa_0$ is held fixed while
$\kappa = \Gamma/N$ varies, then $r = N\kappa_0/\Gamma$ grows with $N$
and the monotonicity strengthens further, so maximum-$N$ selection
remains qualitatively intact.  The marginal case is $\kappa_0 = 0$:
then $r = 0 < \tfrac12$ and Term~1 in~\eqref{eq:H-monotone} is
\emph{negative}, so $H_{N+1} > H_N$ requires $R$ large enough that
Term~2 dominates.  The $r \ge \tfrac{1}{2}$ condition is therefore
specific to the regime with a sufficiently strong central vortex.

The $N = 8$ selection is robust across both limiting assumptions
(fixed~$r$ and fixed~$\kappa_0$) but relies on the Onsager framework
providing the link between $N$, $\Gamma$, and $\kappa_0$, a physical
hypothesis, not a mathematical consequence.
\end{remark}

\noindent
For Jupiter north, $\kappa_0/\kappa \approx 1$ with roughly 30\%
observational uncertainty (Juno; \citealt{Adriani2018}), giving a
conservative range $\kappa_0/\kappa \in [0.7, 1.3]$.  Since
$\kappa_{\mathrm{crit}}(8) = \tfrac{1}{2}$ exactly, this entire
range satisfies the stability condition for $N = 8$.  However,
$\kappa_{\mathrm{crit}}(9) = 1$, which lies \emph{inside} the
uncertainty interval: if $\kappa_0/\kappa \ge 1$, then $N = 9$ is also
Thomson-stable and the Onsager argument would select $N_{\max} = 9$.
The conclusion $N_{\max} = 8$ relies on $\kappa_0/\kappa < 1$, which is
consistent with the central estimate but not guaranteed by current
observational precision.  The observed $N = 8$ is thus evidence that
$\kappa_0/\kappa \lesssim 1$, providing a theory-derived constraint on
the central vortex strength testable with improved Juno wind
measurements.

\subsection{Jupiter's north--south asymmetry}
\label{sec:jupiter-asymmetry}

For Jupiter's south pole ($N = 5$): $\kappa_{\mathrm{crit}}(N) = 0$
for all $N \le 7$, so the Thomson bound alone allows up to $N = 7$
regardless of the central vortex.  The observed $N = 5$ is selected
by the geometric packing constraint
(Proposition~\ref{prop:complementary}, item~3).
South-pole cyclones are individually $\approx 40\%$ larger in
linear size than north-pole cyclones ($r_{\mathrm{core}} \approx
3{,}500$\,km vs.\ $2{,}500$\,km; \citealt{Adriani2018}), and each
cyclone carries a $\beta$-drift exclusion zone extending $\approx
30\%$ beyond the core radius \citep{Gavriel2021}, giving
$r_{\mathrm{excl}} \approx 1.3\,r_{\mathrm{core}} \approx 4{,}550$\,km.
At the south-pole ring radius $R \approx 8{,}000$\,km
($\sim\!7^{\circ}$ from the pole):
\[
  \sin(\pi/5) = 0.588 \ge \frac{r_{\mathrm{excl}}}{R}
  = \frac{4{,}550}{8{,}000} = 0.569,
  \qquad
  \sin(\pi/6) = 0.500 < 0.569.
\]
Thus $N_{\mathrm{pack}} = 5$, which is tighter than the Thomson bound
$N_{\mathrm{Thomson}} = 8$: packing, not stability, is the binding
constraint at Jupiter's south pole.  The selection is tight:
$(\sin(\pi/5) - r_{\mathrm{excl}}/R)/(\sin(\pi/5)) = 3.3\%$.

The entire $N = 8$ vs.\ $N = 5$ asymmetry reduces to a single
dimensionless ratio:
\[
  \chi = \frac{R_d^{(\mathrm{south})}}{R_d^{(\mathrm{north})}}
  \approx \frac{r_{\mathrm{cyclone}}^{(\mathrm{south})}}
              {r_{\mathrm{cyclone}}^{(\mathrm{north})}}
  \approx \frac{3{,}500}{2{,}500} = 1.40.
\]
The constraint intersection requires $\chi \gtrsim 1.23$
for the south-pole packing bound to tighten from $N = 6$ to
$N = 5$.  The observed $\chi = 1.40$ matches to $8\%$.
Deriving $\chi$ from atmospheric physics (insolation asymmetry,
moist convection) remains open: radiative scaling gives
$\chi \sim 1.07$ (insufficient), suggesting that deep convective
dynamics set the pole-to-pole deformation-radius contrast.

\subsection{Saturn's structural stability}
\label{sec:saturn-structural}

The Rossby stationarity argument selects Saturn's wavenumber $n = 6$
(\S\ref{sec:deformation}).
A natural question: can smooth perturbations of the jet destroy the
stationary wave, converting the hexagonal meander into an evanescent
mode?  The answer is no, by a WKB branch-disconnection argument.

The Rayleigh--Kuo equation for stationary ($c = 0$) mode~$n$ on a
zonal jet $U(y)$ reads $\hat\psi'' + V(\rho)\hat\psi = 0$, where
$V = (\beta - U'')/U - n^2/R^2$ is the effective potential in
log-polar coordinates.  At the jet center ($y = 0$) of a Gaussian jet:
\[
  V(0) = \frac{\beta + U_{\max}/\sigma^2}{U_{\max}} - \frac{n^2}{R^2}
  = \frac{\beta_{\mathrm{eff}}}{U_{\max}} - \frac{n^2}{R^2}.
\]
For Saturn ($U_{\max} = 120\;\mathrm{m/s}$, $\sigma = 2.5\times
10^6\;\mathrm{m}$, $\beta = 1.46\times 10^{-12}\;\mathrm{m^{-1}
s^{-1}}$, $n = 6$, $R = R_{\mathrm{S}}\cos 76^\circ = 1.31\times
10^7\;\mathrm{m}$):
$V(0) = 1.72\times 10^{-13} - 2.10\times 10^{-13} =
-3.8\times 10^{-14}\;\mathrm{m}^{-2}$ (negative).
The sign of $V(0)$ depends on $\sigma_{\mathrm{jet}}$:
$V(0) < 0$ for $\sigma_{\mathrm{jet}} \gtrsim 2.3 \times 10^6$\,m,
with a $\sim\!9\%$ margin above the sign change
($(\sigma - \sigma_{\mathrm{crit}})/\sigma_{\mathrm{crit}}$)
at the adopted
$\sigma_{\mathrm{jet}} = 2.5 \times 10^6$\,m (Cassini jet profile
measurements constrain $\sigma_{\mathrm{jet}} \in [2.3, 3.0]
\times 10^6$\,m).

Since $V < 0$, the WKB exponent $s^2 = V < 0$ forces oscillatory
solutions with $s = i\alpha_\rho$ ($\alpha_\rho \ne 0$).  Writing
$s = \sigma + i\alpha_\rho$:
$\operatorname{Im}(s^2) = 2\sigma\alpha_\rho = 0$.
When $\alpha_\rho \ne 0$ (oscillatory region): $\sigma = 0$.
The elliptic ($\sigma = 0$) and hyperbolic ($\alpha_\rho = 0$)
branches are disconnected: smooth jet perturbations cannot move the
solution from the oscillatory branch to the evanescent branch at
leading WKB order.  No continuous deformation of the jet profile
can remove the $n = 6$ stationary wave without first passing through
$V(0) = 0$ (a finite perturbation).  The hexagon is therefore
structurally stable at leading WKB order.

\subsection{Quality metrics and cross-planetary comparison}
\label{sec:quality-metrics}

The two planetary polygon types require different quality metrics.
For \emph{discrete vortex rings} (Jupiter), the natural measure is the
geometric deviation from a perfect $N$-gon:
\begin{equation}\label{eq:sigma-geom}
  \sigma_{\mathrm{geom}} = \frac{1}{N}\sum_{k=1}^{N}
  \bigl|\!\ln|z_k / z_k^{\mathrm{ideal}}|\bigr|,
\end{equation}
where $z_k^{\mathrm{ideal}}$ are the best-fit regular $N$-gon vertices.
For \emph{continuous jet meanders} (Saturn), the appropriate measure is
the fractional radial displacement of the jet center:
\begin{equation}\label{eq:eps-obs}
  \varepsilon_{\mathrm{obs}} = A_n / \bar{R},
\end{equation}
where $A_n$ is the half-amplitude of the $n$-th Fourier mode and
$\bar{R}$ is the mean ring radius.  For Saturn's hexagon
($n = 6$, $\bar{R} \approx 13{,}200\;\mathrm{km}$),
$A_6 \approx 1{,}900\;\mathrm{km}$ and
$\varepsilon_{\mathrm{obs}} \approx 0.14$
\citep{Godfrey1988, Antunano2015}.
These two metrics are not interchangeable: applying
$\sigma_{\mathrm{geom}}$ to Saturn yields $0.50$, an order of magnitude
larger than Jupiter's $0.05$--$0.06$, but this conflates meander
amplitude with positional scatter.

\begin{table}[htbp]
\centering
\caption{Planetary polygon configurations.
$\kappa_0/\kappa$ is estimated from Juno imagery \citep{Adriani2018}
using cyclone area as a proxy ($\sim\!30\%$ systematic uncertainty).
``, '' indicates the quantity is not applicable.}
\label{tab:planets}
\begin{tabular}{lcccl}
\toprule
System & $N$ & Quality metric &
  $\kappa_0/\kappa$ & Selection mechanism \\
\midrule
Saturn hexagon    & 6 & $\varepsilon_{\mathrm{obs}} = 0.14$ & N/A (jet) & Rossby stationarity \\
Jupiter north     & 8 & $\sigma_{\mathrm{geom}} = 0.06$ & $\sim\!1$ & Thomson + Onsager \\
Jupiter south     & 5 & $\sigma_{\mathrm{geom}} = 0.05$ &, & Packing (Prop.~\ref{prop:complementary}) \\
\bottomrule
\end{tabular}
\end{table}

\noindent
Independent numerical support comes from \citet{Siegelman2022}, who
produce $N = 8$ vortex crystals from 2D turbulence simulations
on polar caps, with a central polar cyclone emerging naturally
alongside the ring.  These simulations
include continuous small-scale forcing and large-scale dissipation, yet
the polygon crystal is the persistent stationary state, numerical
evidence that the constrained energy minimum acts as a dynamical
attractor under turbulent forcing within the 2D barotropic framework.


%% ====================================================================
\section{Circulation disorder and stability domain robustness}
\label{sec:disorder-new}

\subsection{Resilience to circulation disorder}
\label{sec:disorder}

In both planetary regimes the vortex circulations $\kappa_k$ are not
perfectly equal; convection injects or removes circulation from individual
vortices continuously.  We model this as $\kappa_k = \kappa(1+\eta_k)$
with i.i.d.\ fractional disorders $\eta_k \sim \mathcal{N}(0,\sigma_\eta^2)$.

We work on $\mathbf{H}^2$ rather than the flat plane because the
hyperbolic setting simultaneously tests two effects: circulation
disorder and curvature stabilization.  In the flat limit
$\rho/a \to 0$ the $\mathbf{H}^2$ eigenvalues recover the Havelock
values exactly, so the flat-plane result is contained as a special
case; conversely, the curvature-stabilized configurations ($N \ge 8$
on $\mathbf{H}^2$) have no flat-plane analog and can only be probed
on the curved surface.

On $\mathbf{H}^2$, the disordered Hamiltonian is
\begin{equation}\label{eq:H_kappa}
  H_\kappa = -\sum_{j<k}\kappa_j\kappa_k\ln|z_j - z_k|
  + \tfrac12\sum_j\kappa_j(K_{\mathrm{tot}}-\kappa_j)\ln(a^2 - |z_j|^2),
\end{equation}
where $K_{\mathrm{tot}} = \sum_k\kappa_k$.  At $\eta_k \equiv 0$ this
reduces exactly to the uniform Hamiltonian $H_{\mathrm{hyp}}$.

Holding the ring positions at the uniform equilibrium (exact to $O(\eta)$;
the $O(\eta)$ shift of equilibrium produces only an $O(\eta^2)$ eigenvalue
correction), we compute the minimum Fourier-mode eigenvalue
$\lambda_{\min}(\eta)$ via finite differences.  Monte Carlo over
$n = 1000$ realizations estimates $P(\lambda_{\min} < 0)$ as a function
of $(\sigma_\eta, N, \rho)$.

\paragraph{Physical range of $\sigma_\eta$.}
Juno imagery \citep{Adriani2018,Gavriel2021} shows Jupiter's eight
north-polar cyclones have peak wind speeds ranging from
$\sim\!60$ to $\sim\!100$\,m\,s$^{-1}$ and diameters from
$\sim\!4{,}000$ to $\sim\!7{,}000$\,km.  Taking circulation
$\kappa \propto V_{\max}\,r_c$, the coefficient of variation across
the eight cyclones is $\sigma_\eta \approx 0.15$--$0.25$.  Saturn's
hexagonal jet is more uniform (Cassini wind profiles give
$\sigma_\eta \lesssim 0.05$ for the azimuthal circulation variation).
The physically relevant range is therefore
$\sigma_\eta \in [0.05, 0.25]$.

\medskip\noindent
Key findings (all uncertainties are $\pm 1\sigma$ binomial standard
errors $\sqrt{P(1{-}P)/n}$ with $n = 1000$):
\begin{itemize}
  \item \textbf{Stable configurations are robust.}  For $N = 6$ at
    $\rho = 0.5a$, the instability probability remains
    below $10\%\;(\pm 0.9\%)$ for $\sigma_\eta \le 0.2$, covering
    the full physically relevant range.

  \item \textbf{Curvature-stabilized rings ($N \ge 8$) are similarly
    robust.}  At $N = 8$, $\rho = 2a$ (well inside the stable region
    on $\mathbf{H}^2$), $P < 0.5\;(\pm 1.6\%)$ for
    $\sigma_\eta \le 0.1$.  At $\sigma_\eta = 0.2$ (Jupiter's
    estimated value), $P$ remains below $15\%$.

  \item \textbf{The $N = 7$ marginal ring is sensitive.}  At $N = 7$,
    $\rho = 0.05a$ (flat-limit regime), $P(\lambda_{\min} < 0)$
    increases monotonically with $\sigma_\eta$:
    $P \approx 0\;(\pm 0.3\%)$ at $\sigma_\eta = 0.1$ and $P > 0$ for
    $\sigma_\eta \ge 0.5$.  The marginal mode ($\lambda_0 = 0$ at
    $K = 0$) is the first to be destabilized by disorder, as expected.
    At Jupiter's estimated $\sigma_\eta \approx 0.2$, a hypothetical
    $N = 7$ ring would have $P \approx 0$, the marginal sensitivity
    manifests only at unphysically large disorder.
\end{itemize}

\noindent
Thus the $N_{\mathrm{crit}}$ boundary is a \emph{stability domain
boundary, not a knife-edge}: configurations comfortably inside the
stable region remain stable under the circulation fluctuations
observed at Jupiter ($\sigma_\eta \approx 0.2$) and Saturn
($\sigma_\eta \lesssim 0.05$).

See Code Availability.

\subsection{The regime criterion}

The dimensionless number $\beta R^3/\kappa$ is the ratio of two
timescales: the Rossby wave period $T_R \sim 1/(\beta R)$ and the
vortex turnover time $T_v \sim R^2/\kappa$.  When
$\beta R^3/\kappa = T_v/T_R \ll 1$, the vortex interactions complete
many orbits before the Rossby wave propagates appreciably, so the
dynamics is vortex-dominated.  When $\beta R^3/\kappa \gg 1$, the
Rossby wave dynamics operates faster than the vortex rearrangement,
so the wave mechanism dominates:
\begin{center}
\begin{tabular}{lcl}
\toprule
Regime & $\beta R^3/\kappa$ & Bridge \\
\midrule
Vortex-dominated & $\ll 1$ & Blob convergence (Proposition~\ref{prop:blob}) \\
Wave-dominated & $\gg 1$ & Rossby stationarity (equation~\eqref{eq:rossby-bridge}) \\
\bottomrule
\end{tabular}
\end{center}

\noindent
\textit{Jupiter's north pole.}
Throughout this regime calculation $\kappa$ denotes the
\emph{Rankine-core} circulation, which integrates only within the
solid-body rotation core and sets the characteristic vortex timescale
$T_v \sim R^2/\kappa$.  The full-cyclone circulation $\kappa_{\mathrm{full}}$
(which includes the outer irrotational skirt) is larger and is the
appropriate quantity for computing orbital periods and eigenfrequencies;
it is used there instead.
The two estimates are numerically distinguished below.

At $84^\circ$N, the planetary vorticity
gradient is $\beta_J = 2\Omega_J\cos(84^\circ)/R_J \approx
5.1\times10^{-13}$\,m$^{-1}$s$^{-1}$ and the cyclone-ring radius
is $R \approx 7{,}470$\,km.  Each Juno-observed circumpolar cyclone has
Rankine core radius $r_c \approx 700$\,km and peak wind $V_{\max}
\approx 80$\,m\,s$^{-1}$ \citep{Adriani2018,Gavriel2021}, giving
$\kappa_{\mathrm{Rank}} = 2\pi r_c V_{\max} \approx 3.5\times10^{8}$\,m$^2$\,s$^{-1}$
(full-cyclone estimate: $r_c \approx 3{,}200$\,km,
$\kappa_{\mathrm{full}} \approx 1.6\times10^9$\,m$^2$\,s$^{-1}$).
Hence
\[
  \frac{\beta R^3}{\kappa}
  = \frac{5.1\times10^{-13}\,(7.47\times10^6)^3}{3.5\times10^8}
  \approx 0.61.
\]
This places Jupiter in the vortex-dominated regime, but only
moderately: $0.61$ is $O(1)$, not $\ll 1$.  Rossby-wave effects
are not negligible at this value, and a decrease in the effective
$\kappa$ (e.g.\ using the Rankine-core rather than full-cyclone
estimate) would push the ratio above unity.  The qualitative
conclusion, discrete cyclones, Thomson stability applies, blob
bridge proves convergence, is robust to this uncertainty because
the vortex-dominated bridge requires only that the cyclones be
identifiable as $N$ distinct vortices (which the Juno imagery
unambiguously confirms), not that $\beta R^3/\kappa$ be
parametrically small.  The ratio enters quantitatively only
through the Rossby-wave correction to the eigenvalues, which is
$O(\beta R^3/\kappa)$ and does not change the sign of
$\lambda_{\min}$ for $N = 8$ (where
$|\lambda_{\min}| = 1 > 0.61$).
Using the full-cyclone circulation
$\kappa_{\mathrm{full}} \approx 1.6\times10^9$\,m$^2$\,s$^{-1}$
gives $\beta R^3/\kappa \approx 0.13$, firmly in the
vortex-dominated regime.

\smallskip\noindent
\textit{Saturn's hexagon.}  At $76^\circ$N, $\beta_S =
2\Omega_S\cos(76^\circ)/R_S \approx 1.3\times10^{-12}$\,m$^{-1}$s$^{-1}$
and $R \approx 14{,}600$\,km.  Saturn's polar jet is
well-described by a Gaussian profile $U = U_{\max}\exp(-y^2/2\sigma^2)$;
its effective circulation is $\kappa = U_{\max}\sqrt{2\pi}\,\sigma$.
With $U_{\max} \approx 100$\,m\,s$^{-1}$ and $\sigma \approx 2{,}100$\,km
(Cassini VIMS wind profiles), $\kappa \approx 5.4\times10^{8}$\,m$^2$\,s$^{-1}$,
giving
\[
  \frac{\beta R^3}{\kappa}
  = \frac{1.3\times10^{-12}\,(1.46\times10^7)^3}{5.4\times10^8}
  \approx 7.6.
\]
Continuous jet meander, Rossby stationarity selects $n = 6$, periodicity
provides quantization.

\smallskip\noindent
Both mechanisms require the 2D logarithmic Green's function as the
common mathematical kernel.


%% ====================================================================
\section{Predictions for quantum systems}
\label{sec:quantum}

\subsection{BEC vortex clusters}

In a quasi-2D Bose--Einstein condensate in a circular trap, quantized
vortices interact logarithmically at scales above the healing length
$\xi$.  The circulation is quantized:
$\kappa = q \cdot h/m$ with integer winding number $q$.
The $\kappa_0/\kappa$ ratio of Proposition~\ref{prop:kappa-crit}
becomes the winding number ratio $q_0/q$, yielding sharp integer-valued
predictions:

\begin{center}
\begin{tabular}{lcl}
\toprule
Configuration & $\kappa_{\mathrm{crit}}$ & Prediction \\
\midrule
$N \le 7$ ring & 0 & Stable without central vortex \\
$N = 8$ ring & $1/2$ & Unstable alone; stable with $q_0 \ge 1$ \\
$N = 9$ ring & $1$ & Marginally stable with $q_0 = 1$;
  stable with $q_0 \ge 2$ \\
$N = 10$ ring & $7/4$ & Unstable with $q_0 = 1$;
  stable with $q_0 \ge 2$ \\
\bottomrule
\end{tabular}
\end{center}

\noindent
The $N = 8$ prediction is the sharpest: a ring of $8$ unit-winding
vortices should be unstable alone and stable upon adding a single
central vortex ($q_0 = 1 > \kappa_{\mathrm{crit}} = 1/2$).

The $N = 7$ case offers a distinct signature.  The quartic restoring
force ($\alpha_0 = 45/14$, Section~\ref{sec:n7}) produces an
\emph{amplitude-dependent} oscillation frequency for the $m = 3$ mode.
In the normal form $\mathcal{H} = (\alpha_0/4)\varepsilon^4$, the
nonlinear oscillation frequency is
$\omega_{m=3} = \alpha_0\,\varepsilon^2\,\kappa/(2\pi R^2)$, where
$\varepsilon$ is the dimensionless perturbation amplitude (radial
displacement / ring radius) and $\kappa = h/m_{\mathrm{atom}}$ is the
single-quantum circulation.  For $^{87}$Rb in a circular trap of
radius $R = 50\;\mu\mathrm{m}$:
$\kappa \approx 1.2\times 10^{-7}\;\mathrm{m^2\,s^{-1}}$,
$\kappa/(2\pi R^2) \approx 7.6\;\mathrm{s^{-1}}$, so
$\omega_{m=3} \approx 24\,\varepsilon^2\;\mathrm{s^{-1}}$.  At a
5\% radial displacement ($\varepsilon = 0.05$):
$\omega_{m=3} \approx 0.06\;\mathrm{s^{-1}}$ (period $\sim\!100\;\mathrm{s}$),
resolvable with standard BEC vortex imaging at $\sim\!10\;\mathrm{Hz}$
frame rates.  The key signature distinguishing the $N = 7$
quartic confinement from the $N \le 6$ quadratic regime is that the
$m = 3$ oscillation frequency \emph{vanishes} as $\varepsilon \to 0$
(rather than approaching a finite value), producing an
amplitude-dependent period that is directly measurable.

The healing length $\xi$ provides a natural blob width
$\varepsilon_{\mathrm{blob}} = \xi/R$.  For typical BEC parameters
($\xi \sim 0.3\;\mu\mathrm{m}$, $R \sim 50\;\mu\mathrm{m}$):
$\varepsilon_{\mathrm{blob}} \approx 0.006$, well within the point
vortex regime.  The blob correction (Proposition~\ref{prop:blob}) is
negligible: $\lambda(0.006) \approx \lambda_0 + O(10^{-4})$.

\emph{Finite-temperature effects.}  Current BEC vortex experiments
operate at finite temperature ($T/T_c \sim 0.3$--$0.7$), where a
thermal cloud coexists with the condensate and exerts a dissipative
drag on vortex motion \citep{Fetter2009}.  This drag damps the
oscillation modes on a timescale $\tau_{\mathrm{drag}} \sim
R^2/(D_v\,T/T_c)$ where $D_v$ is the vortex diffusion coefficient.
For the $N = 7$ quartic mode to be observable, the oscillation period
$(\sim\!100\;\mathrm{s}$ at $\varepsilon = 0.05)$ must be shorter than
$\tau_{\mathrm{drag}}$; at $T/T_c = 0.5$ this requires
$D_v \lesssim 10^{-9}\;\mathrm{m^2\,s^{-1}}$, which is achievable in
large-radius traps ($R \ge 50\;\mu\mathrm{m}$) at low temperature.
The $N = 8$ stability prediction (unstable alone, stable with $q_0 = 1$)
is less sensitive to dissipation because it concerns exponential growth
vs.\ decay, not oscillation frequency; finite temperature adds a damping
channel but does not change the sign of the eigenvalue.

\emph{The hard-wall disk: an exact theorem.}
In a hard-wall circular trap of radius $R$, each
vortex at $z$ has an image at $R^2/\bar{z}$.
The effective interaction is \emph{exactly} the Green's
function on the Poincar\'e disk of radius~$R$
(the conformal equivalence is the identity map).
\begin{proposition}[Hard-wall BEC = Poincar\'e disk]
\label{prop:bec-disk}
In a hard-wall circular BEC of radius~$R$, the
stability eigenvalue of an $N$-gon vortex ring at
radius~$r$ is $\lambda_m = C_1(\xi) - m(N{-}m)/2$
with $\xi = r^2/R^2$ and
$C_1(\xi) = (N{-}1)(1{+}\xi^2)/(1{-}\xi)^2$, the
same formula as for $\mathbf{H}^2$
(Paper~I, equation~\ref{I-eq:C1_H2}).
\end{proposition}
\begin{proof}
The image vortex construction gives the Green's function
$G(z,w) = -(1/2\pi)\log|z{-}w| + (1/2\pi)\log|1{-}z\bar{w}/R^2|$,
which is exactly the Green's function on the Poincar\'e
disk of radius~$R$.
The Havelock decomposition applies without approximation.
\end{proof}

The palindromic thresholds $\xi^*(N)$ are therefore
\emph{exact} predictions for hard-wall BEC experiments:
\begin{center}
\small
\begin{tabular}{cccl}
\toprule
$N$ & $\xi^* = r^2/R^2$ & $r/R$ & Prediction \\
\midrule
$\le 7$ & $0$ & $0$ & Stable at any radius \\
$8$  & $8{-}3\sqrt{7} \approx 0.063$ & $0.25$ &
  Stable for $r > 0.25\,R$ \\
$11$ & $(1{+}\sqrt{2})^{-2} \approx 0.17$ & $0.41$ &
  Stable for $r > 0.41\,R$ \\
$23$ & $\varphi^{-2} \approx 0.38$ & $0.62$ &
  Stable for $r > 0.62\,R$ \\
\bottomrule
\end{tabular}
\end{center}
The \textbf{key testable prediction}: an $8$-vortex ring at
$r/R > 0.25$ in a hard-wall disk BEC should be stable;
at $r/R < 0.25$ it should be unstable to the $m = 4$ mode.
The threshold $r/R = \sqrt{8{-}3\sqrt{7}} \approx 0.25$ is an
algebraic number in $\mathbb{Q}(\sqrt{7})$, the same field
that governs the $7 \to 8$ transition in planetary vortex dynamics.

Flat-bottomed BEC traps, in which the point vortex model has
been experimentally validated
\citep{Gauthier2019}, are the natural setting for testing
these predictions.

\emph{The harmonic trap: exact result.}
For a harmonic trap $V(r) = \tfrac{1}{2}m\omega^2 r^2$
with Thomas--Fermi density $\rho = \rho_0(1 - \xi)$
($\xi = r^2/R_{\mathrm{TF}}^2$), the Fetter precession
analysis \citep{Fetter2001} gives the exact Havelock coefficient
\begin{equation}\label{eq:C1-harmonic}
  C_1^{\mathrm{harm}}(N, \xi)
  = \frac{N-1}{1-\xi},
\end{equation}
a M\"obius function, in contrast to the hard-wall
$C_1^{\mathrm{hard}} = (N{-}1)(1{+}\xi^2)/(1{-}\xi)^2$.
The correction is $O(\xi)$, not $O(\xi^2)$:
$\Delta C_1 = -(N{-}1)\xi(1{+}\xi)/(1{-}\xi)^2$,
a structural change reflecting the soft boundary.
The harmonic-trap threshold solves
$(N{-}1)/(1{-}\xi) = m(N{-}m)/2$, giving the exact
rational values
\begin{equation}\label{eq:xi-harm}
  \xi^*_{\mathrm{harm}} = 1 - \frac{2(N-1)}{m(N-m)},
  \qquad m = \lfloor N/2 \rfloor.
\end{equation}
\begin{center}\small
\begin{tabular}{@{}lccc@{}}
\toprule
 & $\xi^*_{\mathrm{hard}}$ & $\xi^*_{\mathrm{harm}}$
 & $r/R_{\mathrm{TF}}$ (harm) \\
\midrule
$N = 8$  & $8{-}3\sqrt{7} \approx 0.063$ & $1/8 = 0.125$  & $0.354$ \\
$N = 11$ & $0.166$ & $1/3 = 0.333$  & $0.577$ \\
$N = 23$ & $0.385$ & $2/3 = 0.667$  & $0.816$ \\
\bottomrule
\end{tabular}
\end{center}
The harmonic thresholds are all rational, the palindromic
structure and algebraic number fields of the hard-wall case
(\S\ref*{I-sec:palindromic-quad}) do not appear.
The distinction is experimentally accessible: an $N = 8$
ring at $r/R_{\mathrm{TF}} = 0.30$ is predicted stable in
a hard-wall trap but unstable in a harmonic trap.

\emph{Healing-length and quantum corrections.}
The BEC healing length $\xi_{\mathrm{heal}}$ sets the
blob scale $\varepsilon = \sqrt{2}\,\xi_{\mathrm{heal}}/d$
where $d = 2r\sin(\pi/N)$ is the intervortex spacing.
For typical parameters
($\xi_{\mathrm{heal}} \approx 0.3\;\mu\mathrm{m}$,
$R \approx 50\;\mu\mathrm{m}$):
$\varepsilon \approx 0.02$, giving
$\delta\lambda \sim \varepsilon^2 \sim 0.04\%$
(Proposition~\ref{prop:blob}).

The \emph{quantum} correction (from the Bogoliubov--de~Gennes
excitation spectrum) is even smaller.
The BdG anomalous coupling between vortex cores separated
by~$d$ is $O(\exp(-d/\xi_{\mathrm{heal}}))$;
for $N = 8$ at the threshold ($d \approx 9.6\;\mu\mathrm{m}$):
$\exp(-d/\xi_{\mathrm{heal}}) \approx 10^{-14}$.
The effective mass renormalization shifts the precession
frequency but not the energy Hessian (the stability condition
$\lambda_m > 0$ is a property of the potential landscape,
not the kinetic term).
The palindromic threshold $\xi^*(8) = 8{-}3\sqrt{7}$ is
therefore an \emph{exact} feature of the quantum spectrum,
preserved to exponential accuracy by the topological
invariance of the $K_0$ class
(Paper~I, Theorem~\ref{I-thm:index-pairing-k0}).

\emph{Hyperbolic lattice realizations.}
Interacting bosons on a $\{p,q\}$ hyperbolic lattice
(Bose--Hubbard model, hopping~$t$, on-site $U$) reduce
to the point vortex Hamiltonian in the Thomas--Fermi
regime ($U \gg t$):
the mean-field condensate $\psi = \sqrt{\rho}\,e^{i\theta}$
has vortices where $\psi = 0$, and $\theta$ satisfies
$\Delta_{\mathrm{lattice}}\theta = 2\pi\sum_k \delta(z{-}z_k)$,
giving the lattice Green's function interaction
$H = -\sum_{j<k} G_{\mathrm{lattice}}(z_j, z_k)$.
In the continuum limit (lattice spacing $a \ll$ intervortex
spacing~$d$): $G_{\mathrm{lattice}} \to G_{\mathbf{H}^2}$
and the Havelock thresholds apply.
The continuum limit requires $\gtrsim 1000$ lattice
sites, beyond current circuit-QED
($\sim\!100$ sites) but within reach of topoelectric
circuits ($\sim\!10^4$ sites).
On the Bolza tiling ($\{8,8\}$, genus~$2$,
$|\mathrm{Aut}| = 48$): the prediction is
$N_{\mathrm{crit}} = 12$ at the Aut-fixed center,
with $\delta C_1 \approx 0.011$ (Paper~I, \S7).

\subsection{Mesoscopic superconductors}

\begin{remark}[Mesoscopic superconducting disks]
\label{rmk:superconductor}
Abrikosov vortices in a type-II superconducting disk of
radius~$R$ obey the same boundary condition as BEC
vortices in a hard-wall trap (supercurrent vanishes at the
edge $\Leftrightarrow$ Kelvin images at $R^2/\bar{z}$).
Proposition~\ref{prop:bec-disk} therefore applies directly:
$\lambda_m = C_1(\xi) - m(N{-}m)/2$ with $\xi = r_0^2/R^2$.
The observed experimental sequence in mesoscopic Nb disks
(\citealt{Grigorieva2006}):
single polygon for $N \le 6$, marginal polygon or shell
($1{+}6$) at $N = 7$, and shells ($1{+}7$, $2{+}7$, \ldots)
for $N \ge 8$, is consistent with $N_{\mathrm{crit}} = 7$
at the disk center ($\xi = 0$) and central-vortex
stabilization (Section~\ref{sec:central-vortex}) raising
$N_{\mathrm{crit}}$ to~$8$ for the outer ring.
\end{remark}

\subsection{Quantum Hall systems}
The following two remarks apply the stability framework to
fractional quantum Hall states. They are independent of the
gravitational interpretation (Papers~III--IV) and may be
read as standalone applications of the Havelock decomposition.

\begin{remark}[Laughlin quasihole polygons]
\label{rmk:laughlin}
The effective Hamiltonian of the Laughlin state at filling
$\nu = 1/m$ is $H_{\mathrm{eff}} = -m\sum_{j<k}\log|z_j{-}z_k|
+ (4\ell_B^2)^{-1}\sum|z_k|^2$: the logarithmic energy with
harmonic confining.
Quasiholes (created by factors $\prod(z_j{-}\eta_a)$) interact
via $H_{\mathrm{qh}} = -(1/m)\sum_{a<b}\log|\eta_a{-}\eta_b|
+ \text{confining}$, where the $1/m$ comes from the fractional
charge $e^* = e/m$.
For $P$ quasiholes in a ring at radius $r$ inside a Laughlin
droplet of $N$ electrons (radius $R = \ell_B\sqrt{2N}$):
the Havelock eigenvalue is
$\lambda_p = (1/m)\bigl[C_1(\xi) - p(P{-}p)/2\bigr]$
with $\xi = r^2/R^2$.
The factor $1/m$ is overall: \emph{the stability threshold
$P_{\mathrm{crit}}$ is independent of the filling fraction~$m$}.
The palindromic thresholds of Paper~I, \S5 apply directly:
$P \le 7$ is stable at any $r$; $P = 8$ requires
$r/R > 0.25$; etc.
The Berry phase (statistical angle $\theta = \pi/m$, giving
anyonic braiding) adds an imaginary term to $H_{\mathrm{eff}}$
that affects the dynamics but not the energy Hessian, the
stability eigenvalues are real and independent of the
statistical angle.
The conditions:
$N \gtrsim 55$ electrons (for $P = 8$ at $r/R = 0.25$),
and $d = 2r\sin(\pi/P) \gg \ell_B$ (well-separated
quasiholes).
\end{remark}

\begin{remark}[Non-abelian quasiholes and topological quantum computation]
\label{rmk:non-abelian}
The Moore--Read state ($\nu = 5/2$) has quasiholes carrying
Majorana zero modes with non-abelian braiding.
The effective interaction between two quasiholes at
separation~$d$ is
$E(d) = E_{\log}(d) + \delta E_{\mathrm{fusion}}(d)$
where $\delta E_{\mathrm{fusion}} \sim C\exp(-d/\xi_E)$
with $\xi_E \approx 2.3\,\ell_B$
(\citealt{Baraban2009}).
The exponential decay gives effective exponent
$\alpha = \infty \gg 2$ in the RG classification
(Remark~\ref*{I-rem:rg-fixed-point}):
the fusion correction is \emph{deeply irrelevant}.
By the blob bridge (Proposition~\ref{prop:blob}):
$P_{\mathrm{crit}} = 7$ is unchanged.

The same argument applies to all gapped quantum Hall states
(Read--Rezayi, hierarchy states): the fusion correction is
always exponential in $d$ (from the bulk energy gap), hence
always RG-irrelevant.
$P_{\mathrm{crit}} = 7$ is therefore \emph{universal across
all gapped quantum Hall states}, independent of filling
fraction, braiding statistics, and fusion channel.

For topological quantum computation: standard braiding
protocols use $P \le 6$ quasiholes ($4$ for Moore--Read,
$6$ for two-qubit Fibonacci), safely below $P_{\mathrm{crit}}$.
The threshold is not a practical constraint on current
proposals but is a fundamental limit on the complexity of
quasihole polygon configurations in any quantum Hall platform.
\end{remark}

%% Remark on spectral/TQFT independence moved to Paper III (M-013).

%% ====================================================================
\section{Gravitational polygon stability in 2+1 dimensions}
\label{sec:gravity}
%% ====================================================================

The $2{+}1$D gravity literature studies the single
conical deficit \citep{DeserJackiw1984}, two-body scattering,
BTZ black holes, polygon tessellations of spatial slices
(\citealt{tHooft1992}; note that ``polygon'' there refers to
the tiles of the tessellation, not to a configuration of
point masses), and one-loop partition functions
\citep{GiombiMaloneyYin2008}.
What appears not to have been asked is:
\emph{for how many point masses can one form a stable regular
polygon on a spatial slice, and what algebraic number theory
governs the answer?}
This section derives the answer from the vortex stability
framework of Paper~I.

In $2{+}1$ dimensional gravity with negative cosmological
constant ($\Lambda = -1/\ell^2$), spacetime is locally
$\mathrm{AdS}_3$ and has no propagating degrees of freedom.
A point mass $m$ creates a conical deficit of angle
$\alpha = 8\pi G m$ \citep{DeserJackiw1984}.
A constant-time spatial slice is a quotient
$S = \Gamma\backslash\mathbf{H}^2$, and the
\emph{linearized} gravitational potential of a point mass
on~$S$ satisfies $\nabla^2 \Phi = 8\pi G\,m\,\delta(z - z_0)$, the same Laplace equation as for the stream function of a
point vortex.

\begin{proposition}[Gravitational polygon stability]
\label{prop:grav-stability}
In the linearized regime ($8\pi G m_i \ll 2\pi$, small
deficit angles), $N$ equal point masses of mass~$m$ on a
spatial slice $S = \Gamma\backslash\mathbf{H}^2$ interact via
\begin{equation}\label{eq:grav-interaction}
  E_{\mathrm{grav}} \;=\; -(8\pi G)^2\!\sum_{i < j} m^2\,
  G_S(z_i, z_j),
\end{equation}
where $G_S$ is the Green's function of the Laplacian on~$S$
(the same $G_S$ appearing in the vortex energy with
$\kappa \leftrightarrow 8\pi Gm$).
The regular $N$-gon is a critical point of $E_{\mathrm{grav}}$
on the angular-impulse surface, and the Hessian eigenvalues
are
\begin{equation}\label{eq:grav-havelock}
  \lambda_m^{\mathrm{grav}}
  \;=\; C_1(S, \xi) - \frac{m(N{-}m)}{2},
  \qquad m = 1, \ldots, N{-}1,
\end{equation}
with the same $C_1$ and $m(N{-}m)/2$ as for vortices.
The $N$-gon is a stable equilibrium
($\lambda_m^{\mathrm{grav}} > 0$ for all $m$) if and only
if $N \le N_{\mathrm{crit}}(S, \xi)$.
\end{proposition}

\begin{proof}
The linearized gravitational potential satisfies the same
Poisson equation as the vortex stream function, with
$\kappa = 8\pi Gm$.
The interaction energy~\eqref{eq:grav-interaction} is
therefore $E_{\mathrm{grav}} = (8\pi Gm/\kappa)^2
\cdot H_{\mathrm{vortex}}$,
a positive multiple of the vortex Hamiltonian.
At the $N$-gon, the first and second variations of
$E_{\mathrm{grav}}$ are proportional to those of
$H_{\mathrm{vortex}}$, so the critical points, the Hessian
eigenvalues, and the stability criterion
($\lambda_m > 0$ for all $m$) are identical.
The Havelock decomposition (Paper~I,
\S\ref{I-sec:havelock-identity}) gives~\eqref{eq:grav-havelock}.

\emph{Note on dynamics.}
The vortex equations of motion are \emph{first order}
(the Hamiltonian is the stream function), while the
gravitational equations are \emph{second order}
(kinetic $+$ potential energy).
The stability criterion, positive-definite Hessian of the
potential, is the same for both: a positive eigenvalue
$\lambda_m > 0$ gives stable oscillation in the
gravitational case and stable circulation in the vortex case.
A negative eigenvalue gives exponential growth in both.
The criterion $N \le N_{\mathrm{crit}}$ is therefore
identical.
\end{proof}

\emph{Validity and the developing map.}
For $N$ equal masses in a ring, each creates a deficit angle
$\delta = 8\pi Gm$.
The total deficit is $N\delta < 4\pi$ (Gauss--Bonnet for $S^2$),
giving $Gm < 1/(2N)$.
We now show that the backreaction of the deficit
angles on the flat-plane stability is \emph{exactly zero}.

The \textbf{developing map} of the flat cone metric with
$\mathbb{Z}_N$ symmetry is
\begin{equation}
  \label{eq:developing-map}
  f(z) \;=\; z \cdot {}_2F_1\!\bigl(\beta,\, \tfrac{1}{N};\,
    1 {+} \tfrac{1}{N};\, z^N\bigr),
  \qquad \beta = 4Gm,
\end{equation}
with derivative $f'(z) = (1 - z^N)^{-\beta}$.
This is a conformal isometry from the cone surface to the
Euclidean plane: the metric is $|f'|^2|dz|^2$, flat away
from the cone tips, with cone angle $2\pi(1{-}\beta)$ at
each $N$-th root of unity.

\emph{$\mathbb{Z}_N$ symmetry.}
$f(\omega z) = \omega f(z)$ for $\omega = e^{2\pi i/N}$:
the image of the $N$-gon of cone points is a regular $N$-gon
in the developed plane, with radius
\begin{equation}
  \label{eq:image-radius}
  R_f \;=\;
  \frac{\Gamma(1 + 1/N)\,\Gamma(1-\beta)}{\Gamma(1 + 1/N - \beta)}
\end{equation}
(Gauss's summation formula for ${}_2F_1$ at $z = 1$).
At $\beta = 0$: $R_f = 1$.
At the Gauss--Bonnet endpoint $\beta = 1/N$:
$R_f = (\pi/N)/\!\sin(\pi/N)$ (by the $\Gamma$-reflection
formula).
The image radius \emph{increases} with $\beta$: the cone metric
inflates distances near the tips.

\begin{proposition}[Conformal invariance of flat-plane stability]
\label{prop:conformal-invariance}
On a flat spatial slice ($\Lambda = 0$) with $N$ equal conical
singularities of deficit $\delta = 8\pi Gm$ in a regular ring,
the Havelock eigenvalues are
\[
  \lambda_m \;=\; (N{-}1) - \tfrac{1}{2}m(N{-}m)
\]
for all $Gm \in [0, 1/(2N))$.
In particular, $N_{\mathrm{crit}} = 7$ is independent of the
deficit angle.
\end{proposition}

\begin{proof}
The two-dimensional Green's function is conformally invariant:
for any conformal metric $g = e^{2u}|dz|^2$ on the plane,
$\Delta_g G = \delta$ gives
$G(z,w) = -(2\pi)^{-1}\log|z{-}w|$
in conformal coordinates, independent of $u$.
Therefore the interaction energy
$H = -\sum_{j<k}\log|z_j - z_k|$ is the same in conformal
coordinates as on the flat plane.

The $\mathbb{Z}_N$ symmetry maps the $N$-gon to a regular
$N$-gon of radius $R_f$ in the developed plane.
The angular interaction
$-\sum_{j<k}\log|e^{i\theta_j} - e^{i\theta_k}|$
is independent of the radius (since $\log R$ drops from the
Hessian), and so are the Fourier eigenvalues
$m(N{-}m)/2$ of the tangential Hessian.
The constraint curvature contributes $+(N{-}1)$
independently of $R_f$, completing the
Havelock decomposition.
\end{proof}

\emph{The topology change at $\beta = 1/N$.}
At the Gauss--Bonnet endpoint $Gm = 1/(2N)$, the total deficit
exhausts the full curvature budget and the spatial slice
\emph{compactifies} to $S^2$.
On the sphere, the Green's function acquires a
curvature correction:
$G^{S^2} = -(2\pi)^{-1}\log|z{-}w|
+ (4\pi)^{-1}\log[(1{+}|z|^2)(1{+}|w|^2)]$
(in stereographic coordinates).
The additional term replaces $C_1 = N{-}1$ with
$C_1 = (N{-}1)\cos\varphi < N{-}1$, giving
$N_{\mathrm{crit}}^{S^2} \le 6$.
The transition from $N_{\mathrm{crit}} = 7$ (flat, all
$Gm < 1/(2N)$) to $N_{\mathrm{crit}} \le 6$ (sphere) is
therefore \emph{discontinuous}: it occurs at the topology
change, not through a gradual backreaction.
The curvature-stability dichotomy of
Proposition~\ref{prop:grav-ncrit} arises from the cosmological
constant $\Lambda$ (which sets $\xi = R^2/\ell^2$), not from
the deficit angles.

\emph{A remark on ``no gravitational force in $2{+}1$D.''}
Pure $2{+}1$D gravity has no propagating degrees of freedom
and no Newtonian $1/r$ force at a distance.
But point masses on an AdS$_3$ spatial slice
$\Gamma\backslash\mathbf{H}^2$ \emph{do} have a
well-defined static potential: the Green's function of the
Laplacian, which is logarithmic.
This potential does not propagate dynamically (there are no
gravitational waves), but it determines the equilibrium
energy landscape.
The stability analysis depends only on the second variation
of this static energy, for which the Havelock decomposition
applies exactly as in the vortex case.

Proposition~\ref{prop:grav-stability} transfers the full
palindromic hierarchy to gravity:
\begin{itemize}
  \item $N_{\mathrm{crit}} = 7$ on any spatial slice
    (universality, Paper~I, Theorem~\ref{I-thm:universality}).
  \item The palindromic thresholds of Paper~I, \S6 are
    \emph{gravitational phase transitions}: algebraic-number
    curvature values at which the stability of the mass
    polygon changes.
  \item The five algebraic-integer thresholds
    $(\xi^*(8) = 8{-}3\sqrt{7}$, \ldots,
    $\xi^*(23) = \varphi^{-2})$ give parameter-free
    predictions for the stability of mass rings in $2{+}1$D
    AdS.
\end{itemize}

%% ====================================================================
\subsection{The gravitational prediction}
\label{sec:grav-prediction}

The results of this section yield a concrete, parameter-free
prediction for $2{+}1$D gravity.

\begin{proposition}[Maximum stable polygon in $2{+}1$D gravity]
\label{prop:grav-ncrit}
In linearized $2{+}1$D gravity with cosmological constant
$\Lambda$ ($Gm \ll 1/(4N)$, weak-field regime), the maximum
number of equal-mass conical singularities that can form a
stable regular polygon on a spatial slice satisfies
\[
  \mathrm{sgn}(N_{\mathrm{crit}} - 7)
  \;=\; -\mathrm{sgn}(\Lambda).
\]
Specifically: $N_{\mathrm{crit}} = 7$ for $\Lambda = 0$;
$N_{\mathrm{crit}} \le 6$ for $\Lambda > 0$; and
$N_{\mathrm{crit}} \ge 8$ for $\Lambda < 0$, with
$N_{\mathrm{crit}}$ non-decreasing in $|{\Lambda}|$
(Proposition~\ref{I-prop:ncrit-monotone}).
The transitions occur at the palindromic thresholds
$\xi^*(N)$ of Paper~I, \S5, with
$\xi = R^2/\ell_{\mathrm{AdS}}^2$.
\end{proposition}

\begin{proof}
The gravitational Havelock decomposition
(Proposition~\ref{prop:grav-stability}) gives
$\lambda_m = C_1(\xi) - m(N{-}m)/2$.
The stability condition $\lambda_m > 0$ for all $m$,
the monotonicity of $C_1(\xi)$, and the curvature-stability
dichotomy (Paper~I,
Proposition~\ref{I-prop:curvature-dichotomy}) transfer
directly.
\end{proof}

\paragraph{Validity.}
The weak-field condition and the stability condition are
\emph{independent}: $Gm$ controls the deficit angle
(and hence the validity of the logarithmic approximation),
while $\xi = R^2/\ell^2$ controls the ring position
(and hence the stability).
At $1\%$ of the maximum mass ($Gm = 0.01/(4N)$): the total
deficit angle is $\approx 3.6^\circ$ for all~$N$, well within
the logarithmic regime.
The constraint structures (angular momentum in gravity,
angular impulse in vortex dynamics) are isomorphic to
$O(Gm)$: both arise from the SO$(2)$ rotational symmetry
of the polygon, and the relativistic corrections to the
gravitational angular momentum shift the Havelock eigenvalue
by $O(Gm) \sim 1\%$, negligible against binding eigenvalues
of $O(1)$.


%% Three-layer decomposition, orbifold--Havelock correspondence,
%% equivariant Riemann--Roch, and WDW bridge moved to Paper III (M-012).
%% See Paper~III, \S\ref{III-sec:three-layer} and \S\ref{III-sec:wdw}.

\paragraph{Code availability.}
All numerical verifications are in the companion repository
(\url{https://github.com/gspea/planetary-polygons-unified}).
Parameter accounting and predictions appear in Paper~VII.

\appendix
\section{Infinite-dimensional convergence proof}
\label{app:convergence}

The full vorticity perturbation space is
infinite-dimensional, including blob \emph{shape} deformations.
The Hessian of the 2D Euler energy
$E[\omega] = \tfrac12\iint\omega(x)G(x,y)\omega(y)\,\dd x\,\dd y$ at
$\omega_\varepsilon = \sum_k \kappa\,\phi_\varepsilon(\cdot - z_k)$
decomposes into three sectors:

\emph{(a) Center translations.}  The Hessian restricted to the
$2N$-dimensional subspace of center displacements is \emph{exactly}
the blob Hessian of Proposition~\ref{prop:blob}.  Convergence to the
point vortex Hessian at $O(\varepsilon^\alpha)$ follows from the
convolution estimates.

\emph{(b) Shape perturbations.}  Mean-zero, center-preserving
deformations enter through the self-energy.  The second variation at
mean-zero $\delta\omega$ is:
\begin{equation}\label{eq:self-energy-Hminus1}
  \delta^2 E_{\mathrm{self}}
  = \frac{1}{4\pi^2}\,\|\delta\omega\|_{H^{-1}}^2.
\end{equation}
The rescaling $\phi_\varepsilon(x)
= \varepsilon^{-2}\phi(x/\varepsilon)$ gives
$E_{\mathrm{self}}[\phi_\varepsilon] = c/\varepsilon^2$,
so shape modes are \emph{unconditionally stable}.

\emph{(c) Center-shape coupling.}  The vanishing monopole and dipole
moments of $\delta\psi$ give quadrupole coupling:
$C_{jk} = O(\varepsilon^2/d_{\min}^3)$
(numerically verified for the Cauchy blob at $d_{\min} \ge 3\varepsilon$).
The ratio to the shape eigenvalue
is $O(\varepsilon^4/d_{\min}^3) \to 0$ as $\varepsilon\to 0$.

\begin{proposition}[Inertia preservation]
\label{prop:inertia}
For $\varepsilon < \varepsilon_0(N)$, where
$\varepsilon_0 = (\lambda_{\min}\,d_{\min}^6)^{1/4}$,
the full Hessian of
$E[\omega_\varepsilon] + \mu_L\,L[\omega_\varepsilon]$ on the space of
$\mathbb{Z}_N$-symmetric vorticity perturbations has the same inertia
as the point vortex Lagrangian Hessian of
Theorem~\ref*{I-thm:constrained-min}.
\end{proposition}

\begin{proof}
Write the constrained full Hessian in block form
$\mathcal{H} = \bigl(\begin{smallmatrix} H_{cc} & H_{cs} \\ H_{cs}^\top & H_{ss} \end{smallmatrix}\bigr)$.
By~\eqref{eq:self-energy-Hminus1},
$\delta^2 E_{\mathrm{self}} = (1/4\pi^2)\|\delta\omega\|_{H^{-1}}^2$.
On the disk of radius~$\varepsilon$, the Poincar\'e inequality
$\|\delta\omega\|_{H^{-1}} \ge (c_P/\varepsilon)\|\delta\omega\|_{L^2}$
for mean-zero perturbations gives
$H_{ss} \ge (c_P^2/(4\pi^2\varepsilon^2))\,\mathrm{Id}$, hence
$H_{ss}$ is coercive.
By Proposition~\ref{prop:blob},
$\lambda_c^{(k)}(\varepsilon) = \lambda_k^{\mathrm{Hav}} + O(\varepsilon^2)$.
From~(c), $\|H_{cs}\|_{\mathrm{op}} = O(\varepsilon^2/d_{\min}^3)$.
By the Schur complement, the correction is
$O(\varepsilon^6/d_{\min}^6) \to 0$, preserving signs.
\end{proof}

The explicit concentration condition is
\begin{equation}\label{eq:concentration-bound}
  \frac{\varepsilon}{d_{\min}} < \delta(N)
  := \sqrt{\frac{\lambda_{\min}(N)}{6(2N{-}3)}},
\end{equation}
giving $\delta(4) = 0.18$, $\delta(5) = 0.15$, $\delta(6) = 0.096$.

\begin{corollary}[Inertia preservation for $N \le 6$]
\label{cor:inertia-N6}
For $N \le 6$ and $\varepsilon/d_{\min} < \delta(N)$,
the full Hessian inertia equals the Havelock inertia.
\end{corollary}

\begin{corollary}[Marginal-mode resolution for $N = 7$]
\label{cor:inertia-N7}
For $N = 7$, $\delta(7) = 0$, but
$\lambda_{\min}(\varepsilon) > 0$ for all $\varepsilon \in (0, d_{\min})$:
in Region~1 ($\varepsilon < 2/\sqrt{|c_7'|}$),
$\lambda \ge \varepsilon^2(4 - |c_7'|\varepsilon^2) > 0$;
in Region~2, interval arithmetic on a $2{,}000$-point grid
gives $\lambda \ge 0.34 > 0$.
\end{corollary}

\bibliographystyle{plainnat}
\bibliography{../shared/refs}
\end{document}
